%-%-%-%-%-%-%-%-%-%-%-%-%-%-%-%-%-%-%-%-%-%-%-%-%-%-%-%-%-%-%-%-%-%-%-%-%-%-%
%%% MAIN DOCUMENT %%%

%-%-%-%-%-%-%-%-%-%-%-%-%-%-%-%-%-%-%-%-%-%-%-%-%-%-%-%-%-%-%-%-%-%-%-%-%-%-%

\documentclass[12pt, oneside]{book}


%%% AESTHETICS %%%
%-%-%-%-%-%-%-%-%-%-%-%-%-%-%-%-%-%-%-%-%-%-%-%-%-%-%-%-%-%-%-%-%-%-%-%-%-%-%


%%% Dimensions and Spacing %%%
\usepackage[margin=1in]{geometry}
\usepackage{setspace}
\linespread{1}

%%% Define new colors %%%
\usepackage{xcolor}
\definecolor{orangehdx}{rgb}{0.96, 0.51, 0.16}

% Normal colors
\definecolor{xred}{HTML}{BD4242}
\definecolor{xblue}{HTML}{4268BD}
\definecolor{xgreen}{HTML}{52B256}
\definecolor{xpurple}{HTML}{7F52B2}
\definecolor{xorange}{HTML}{FD9337}
\definecolor{xdotted}{HTML}{999999}
\definecolor{xgray}{HTML}{777777}
\definecolor{xcyan}{HTML}{80F5DC}
\definecolor{xpink}{HTML}{F690EA}
\definecolor{xgrayblue}{HTML}{49B095}
\definecolor{xgraycyan}{HTML}{5AA1B9}

% Dark colors
\colorlet{xdarkred}{red!85!black}
\colorlet{xdarkblue}{xblue!85!black}
\colorlet{xdarkgreen}{xgreen!85!black}
\colorlet{xdarkpurple}{xpurple!85!black}
\colorlet{xdarkorange}{xorange!85!black}
\definecolor{xdarkcyan}{HTML}{008B8B}
\colorlet{xdarkgray}{xgray!85!black}

% Very dark colors
\colorlet{xverydarkblue}{xblue!50!black}

% Document-specific colors
\colorlet{normaltextcolor}{black}
\colorlet{figtextcolor}{xblue}

% Light colors
\colorlet{xlightblue}{xblue!85!white}

% Enumerated colors
\colorlet{xcol0}{black}
\colorlet{xcol1}{xred}
\colorlet{xcol2}{xblue}
\colorlet{xcol3}{xgreen}
\colorlet{xcol4}{xpurple}
\colorlet{xcol5}{xorange}
\colorlet{xcol6}{xcyan}
\colorlet{xcol7}{xpink!75!black}

% Blue-Purple (should just used colorbrewer...)
\definecolor{xrainbow0}{HTML}{e41a1c}
\definecolor{xrainbow1}{HTML}{a24057}
\definecolor{xrainbow2}{HTML}{606692}
\definecolor{xrainbow3}{HTML}{3a85a8}
\definecolor{xrainbow4}{HTML}{42977e}
\definecolor{xrainbow5}{HTML}{4aaa54}
\definecolor{xrainbow6}{HTML}{629363}
\definecolor{xrainbow7}{HTML}{7e6e85}
\definecolor{xrainbow8}{HTML}{9c509b}
\definecolor{xrainbow9}{HTML}{c4625d}
\definecolor{xrainbow10}{HTML}{eb751f}
\definecolor{xrainbow11}{HTML}{ff9709}

%------- %
% XHFILL %
%------- %


%%% Chapter Headings %%%
\usepackage[Glenn]{fncychap}
\ChTitleVar{\bfseries\scshape\color{black}}
\ChNumVar{\large\selectfont\color{black}}
\ChNameVar{\large\color{black}}
\usepackage{xpatch}
\xpatchcmd\DOCH
{\mghrulefill}{\color{orangehdx}\mghrulefill}
{}{\PatchFailed}
\xpatchcmd\DOTI
{\mghrulefill}{\color{orangehdx}\mghrulefill}
{}{\PatchFailed}
\xpatchcmd\DOTIS
{\mghrulefill}{\color{orangehdx}\mghrulefill}
{}{\PatchFailed}

%% Change Chapter Heading Placement %%
\usepackage{etoolbox}
\makeatletter
\patchcmd{\@makechapterhead}{\vspace*{50\p@}}{\vspace*{-20\p@}}{}{}
\patchcmd{\@makeschapterhead}{\vspace*{50\p@}}{\vspace*{-20\p@}}{}{}
\patchcmd{\DOTI}{\vskip 80\p@}{\vskip 40\p@}{}{}
\patchcmd{\DOTIS}{\vskip 40\p@}{\vskip 0\p@}{}{}
\makeatother

\renewcommand{\thesection}{\thechapter.\arabic{section}} %% Chapter.Section Numbering


%%% FIGURES %%%
\usepackage{graphicx}  
\graphicspath{ {images/} }  
% \numberwithin{figure}{section}
\usepackage{float}
\usepackage{caption}

%%% Hyperlinks %%%
\usepackage{hyperref}
\definecolor{horange}{HTML}{f58026}
\hypersetup{
	colorlinks=true,
	linkcolor=horange,
	filecolor=horange,      
	urlcolor=horange,
}

%% Headers and Footers %%
\usepackage{fancyhdr} % This should be set AFTER setting up the page geometry
\pagestyle{fancy} % options: empty , plain , fancy
\renewcommand{\chaptermark}[1]{\markboth{\thechapter.\ #1}{}}
\fancyhead[R]{\leftmark}
\fancyhead[L]{Hendrix College}
\fancyhead[C]{\includegraphics[height=.50cm]{images/small logo.png}}
\usepackage{xpatch}
\xpretocmd\headrule{\color{orangehdx}}{}{\PatchFailed}
\setlength{\footskip}{0.5in}
\setlength{\headheight}{15pt} % Adjust accordingly if you have a tall header


%%%%% Colored Boxes %%%%%
\usepackage{tcolorbox}
% \numberwithin{BoxCounter}{section}


\newtcolorbox{note}{colback=black!15!white, colframe=black, boxrule=0.5mm, arc=4mm, left=4mm, right=4mm, top=2mm, bottom=2mm}

%% Theorems %%
\newtcolorbox[]{theorem}[2][]{%
    rounded corners,
    colframe=xblue, %color of frame
    colback=xblue!15, %color of background
    coltitle=white, %color of title text
    label={thm:#2},
    title={\large \textbf{Theorem: #2}},
    #1
}

% %% This command needs to exist for the square root of 2 theorem to work

% \newtcolorbox[]{sqrtTh}[2][]{%
%     rounded corners,
%     colframe=xblue, %color of frame
%     colback=xblue!15, %color of background
%     coltitle=white, %color of title text
%     label={thm:#2},
%     title={\large \textbf{Theorem: \texorpdfstring{The Irrationality of $\sqrt{2}$}{The Irrationality of sqrt(2)}}},
%     #1
% }

%% Definitions %%
\newenvironment{definition}[2][]{%
    \textbf{Definition \Gls{#2}:} #1
    \par\vspace{0.5em}\noindent
}{\par\vspace{1em}\noindent}

\newcommand{\wrpdef}[2]{%
    \newglossaryentry{#1}{name=#1,description={#2}}
    \begin{definition}
        {#1}#2.
    \end{definition}
}

%% Axioms %%
\newtcolorbox[]{axiom}[2][]{%
    rounded corners,
    colframe=xdarkgreen,
    colback=xdarkgreen!15,
    coltitle=white,
    label={ax:#2},
    title={\large \textbf{Axiom of #2}},
    #1
}

%% Exercises %%
\newtcolorbox[]{exercise}[2][]{%
    rounded corners,
    colframe=orangehdx!75,
    colback=orangehdx!10,
    coltitle=white,
    label={ex:#2},
    title={\large \textbf{Exercise #2}},
    before={\setstretch{1}}, % Set line spacing to 1 before the box content
    after={\setstretch{1}},  % Reset line spacing after the box content
    #1,
}


%% Examples %%
% need to add counter here: use counter=boxcounter
\newtcolorbox[]{example}[2][]{%
    rounded corners,
    colframe=teal,
    colback=teal!15,
    coltitle=white,
    label={ex:},
    title={\large \textbf{Example: #2}},
    #1,
}

%% Corollaries %%
\newtcolorbox[]{corollary}[3][]{%
    rounded corners,
    colframe=xpurple,  
    colback=xpurple!15,        
    coltitle=white,
    title={\large \textbf{Corollary to #3: #2}},
    #1
}

%-%-%-%-%-%-%-%-%-%-%-%-%-%-%-%-%-%-%-%-%-%-%-%-%-%-%-%-%-%-%-%-%-%-%-%-%-%-%
%% MATH PACKAGES, ENVIRONMENTS, COMMANDS %%
%-%-%-%-%-%-%-%-%-%-%-%-%-%-%-%-%-%-%-%-%-%-%-%-%-%-%-%-%-%-%-%-%-%-%-%-%-%-%
%You'll need your own packages, theroem types, and commands.

\usepackage{tikz}
\usepackage{pifont}
\usetikzlibrary{shapes,backgrounds,calc,patterns}
\usepackage{amsmath,amsthm} 
\usepackage{mathtools}
\usepackage{amssymb}
\usepackage{mdframed}
\usepackage{tabularx}
\usepackage{mathrsfs}
\usepackage{multicol}
\usepackage{slashed}
\usepackage{enumerate}
\usepackage{enumitem}
\usepackage{lipsum}  %This package lets us generate random text for example purposes.
\usepackage{glossaries}

\makeglossaries



%%% New Environments %%%
\newtheorem{exmp}{Example}[section]

%%% Custom Comands %%%

\newcommand{\cmark}{\ding{51}}%
\newcommand{\xmark}{\ding{55}}%

% Natural Numbers 
\newcommand{\N}{\ensuremath{\mathbb{N}}}

% Whole Numbers
\newcommand{\W}{\ensuremath{\mathbb{W}}}

% Integers
\newcommand{\Z}{\ensuremath{\mathbb{Z}}}

% Rational Numbers
\newcommand{\Q}{\ensuremath{\mathbb{Q}}}

% Real Numbers
\newcommand{\R}{\ensuremath{\mathbb{R}}}

% Complex Numbers
\newcommand{\C}{\ensuremath{\mathbb{C}}}

\newcommand{\pfs}{\noindent\makebox[\linewidth]{\rule{\textwidth}{0.4pt}}\vspace{0.5cm}}

\newcommand{\caution}{\marginnote{\includegraphics[width=2em]{caution.png}}}

\newmdenv[
  topline=false,
  bottomline=false,
  rightline=false,
  leftline=true,
  linewidth=1.5pt,
  linecolor=black, % default color, will be overridden in custom commands
  innertopmargin=0pt,
  innerbottommargin=0pt,
  innerrightmargin=10pt,
  innerleftmargin=10pt,
  leftmargin=0pt,
  rightmargin=0pt,
  skipabove=\topsep,
  skipbelow=\topsep,
]{customframedproof}

\newenvironment{proofpart}[2][black]{
    \begin{mdframed}[
        topline=false,
        bottomline=false,
        rightline=false,
        leftline=true,
        linewidth=1pt,
        linecolor=#1!40, % Custom color
        innertopmargin=10pt,
        innerbottommargin=10pt,
        innerleftmargin=10pt,
        innerrightmargin=10pt,
        leftmargin=0pt,
        rightmargin=0pt,
        skipabove=\topsep,
        skipbelow=\topsep
    ]
    \noindent #2
}{
    \end{mdframed}
}

\newcommand{\sbseteqpf}[4]{%
    \begin{customframedproof}[linecolor=#4]
        \begin{proof}
                We need to show these expressions are subsets of each other so we can utilize \thm{subset} to prove they are equal: 
            \begin{proofpart}[#4]{(\(\subseteq\))}
                #1
            \end{proofpart}
    \vspace{2mm}
            \begin{proofpart}[#4]{(\(\supseteq\))}
                #2
            \end{proofpart}
    \vspace{2mm}
            Therefore, by \thm{subset}, #3
        \end{proof}
    \end{customframedproof}
}

\newcommand{\iffpf}[4]{
    \begin{customframedproof}[linecolor=#4]
        \begin{proof}
            We show this by proving both implications:
            \begin{proofpart}[#4]{(\(\Rightarrow\))}
                #1
            \end{proofpart}
    \vspace{2mm}
            \begin{proofpart}[#4]{(\(\Leftarrow\))}
                #2
            \end{proofpart}
    \vspace{2mm}
            #3
        \end{proof}
    \end{customframedproof}
}

% Command for lemma proofs with a predefined color
\newcommand{\expf}[1]{
    \begin{customframedproof}[linecolor=orangehdx!75]
        \begin{proof}
        #1
        \end{proof}
    \end{customframedproof}
}


% Command for theorems proofs with a predefined color
\newcommand{\tpf}[1]{
    \begin{customframedproof}[linecolor=xblue]
        \begin{proof}
        #1
        \end{proof}
    \end{customframedproof}
}
% Command for corollaries proofs with a predefined color
\newcommand{\cpf}[1]{
    \begin{customframedproof}[linecolor=purplehdx]
        \begin{proof}
        #1
        \end{proof}
    \end{customframedproof}
}

% Need to add in customization for examples
\newcommand{\sol}[1]{
    \begin{customframedproof}[linecolor=orangehdx!75]
        \begin{solution}
        #1
        \end{solution}
    \end{customframedproof}
}

\newcommand{\pb}[2]{
    \begin{exercise}{#1}[label=ex:#1]
    #2
    \end{exercise}
}

\def \proofDistance {10pt}
\newenvironment{solution}
  {\textit{Solution.}}

% Command for theorems proofs with a predefined color
\newcommand{\epf}[1]{
    \begin{customframedproof}[linecolor=teal]
        \begin{proof}
        #1
        \end{proof}
    \end{customframedproof}
}

\renewcommand{\theenumi}{\alph{enumi}} 
\renewcommand{\labelenumi}{(\theenumi)}

%-%-%-%-%-%-%-%-%-%-%-%-%-%-%-%-%-%-%-%-%-%-%-%-%-%-%-%-%-%-%-%-%-%-%-%-%-%-%
   %% CODE FORMATING %%%
%-%-%-%-%-%-%-%-%-%-%-%-%-%-%-%-%-%-%-%-%-%-%-%-%-%-%-%-%-%-%-%-%-%-%-%-%-%-%

\usepackage{listings}
\lstdefinestyle{javaStyle}{
    language=Java,
    basicstyle=\ttfamily\small,
    keywordstyle=\bfseries\color{blue},
    commentstyle=\color{green},
    stringstyle=\color{red},
    showstringspaces=false,
    breaklines=true,
    numbers=left,
    numberstyle=\tiny\color{gray},
    stepnumber=1,
    numbersep=5pt,
    captionpos=b,
    frame=single,
    tabsize=4,
    upquote=true
}


%-%-%-%-%-%-%-%-%-%-%-%-%-%-%-%-%-%-%-%-%-%-%-%-%-%-%-%-%-%-%-%-%-%-%-%-%-%-%




% end of preamble
%-%-%-%-%-%-%-%-%-%-%-%-%-%-%-%-%-%-%-%-%-%-%-%-%-%-%-%-%-%-%-%-%-%-%-%-%-%-%

\begin{document}
	
%-%-%-%-%-%-%-%-%-%-%-%-%-%-%-%-%-%-%-%-%-%-%-%-%-%-%-%-%-%-%-%-%-%-%-%-%-%-%
 %%% COVER PAGE %%%
%-%-%-%-%-%-%-%-%-%-%-%-%-%-%-%-%-%-%-%-%-%-%-%-%-%-%-%-%-%-%-%-%-%-%-%-%-%-%

% Do not use all caps.
\newcommand{\titlestandin}[0]{Mathematical Cryptography}
\newcommand{\cussubtitle}[0]{MATH 490}
% Date format should be like: "January 1, 2001"
\newcommand{\startdate}[0]{August 26, 2024}
\newcommand{\customenddate}[0]{December 2, 2024}
% Be sure to include degree recognition (e.g., B.S., M.S., Ph.D)
\newcommand{\professor}[0]{Prof. Allie Ray, Ph.D.}

%-%-%-%-%-%-%-%-%-%-%-%-%-%-%-%-%-%-%-%-%-%-%-%-%-%-%-%-%



\begin{titlepage}
\begin{center}

\vspace*{-2cm}
\includegraphics[width=0.8\textwidth]{images/Hendrix Logo.png}\\
\vfill

% Horizontal line above the title in 'horange' color
\textcolor{horange}{\rule{\textwidth}{1.0pt}}

\vspace{2em}

{\huge \textbf{\titlestandin}}

\vspace{1em} % Space between the title and the bottom line

\textcolor{horange}{\rule{\textwidth}{1.0pt}}

\vspace*{1\baselineskip}

{\LARGE \textbf{\cussubtitle}}

\begin{large}
\vspace*{2\baselineskip}

\textit{Start}  \\[1ex]
{\scshape \startdate} \\[0.3\baselineskip] % Year published

\vspace*{1\baselineskip}

\emph{Author} \\[1ex]
%Submitted by \\[\baselineskip]
{\Large Paul Beggs \\ \par} % Editor list
{\href{mailto:BeggsPA@Hendrix.edu}{{BeggsPA@Hendrix.edu}}}\\ % Editor affiliation

\vspace*{1\baselineskip}

\textit{Instructor} \\[1ex] % Tagline(s) or further description
\professor

\vspace*{1\baselineskip}

\textit{End}\\[1ex]
{\scshape  \customenddate} \\[0.3\baselineskip] % Year published

\thispagestyle{empty} 

\end{large}
\end{center}
\end{titlepage}
\pagebreak

%-%-%-%-%-%-%-%-%-%-%-%-%-%-%-%-%-%-%-%-%-%-%-%-%-%-%-%-%-%-%-%-%-%-%-%-%-%-%

%%% Table of Contents %%%

% No need to edit.

%-%-%-%-%-%-%-%-%-%-%-%-%-%-%-%-%-%-%-%-%-%-%-%-%-%-%-%-%-%-%-%-%-%-%-%-%-%-%
	
	% \thispagestyle{plain}
	% % Table of Contents will automatically update based on included chapters
	% % You may need to compile twice for the toc to update
	% \begin{spacing}{1} %Can change spacing between entries in toc
	% 	\renewcommand{\contentsname}{\Large\textbf{Table of Contents}} %Can rename here
	% 	\tableofcontents
	% 	\addtocontents{toc}{\protect\enlargethispage{\baselineskip}} 
	% \end{spacing}


%-%-%-%-%-%-%-%-%-%-%-%-%-%-%-%-%-%-%-%-%-%-%-%-%-%-%-%-%-%-%-%-%-%-%-%-%-%-%



%-%-%-%-%-%-%-%-%-%-%-%-%-%-%-%-%-%-%-%-%-%-%-%-%-%-%-%-%-%-%-%-%-%-%-%-%-%-%



%%% CHAPTERS %%%

% The document will automatically update (including numbering and toc)
% You may need to compile twice to update the toc
% Witness magic!
	
% Chapter files -- containing the actual text of the document
 % These are in the chapters folder
% All images (for any chapter) must be placed in the images folder

% Example chapters are included with examples on formatting boxes and figures


%-%-%-%-%-%-%-%-%-%-%-%-%-%-%-%-%-%-%-%-%-%-%-%-%-%-%-%-%-%-%-%-%-%-%-%-%-%-%

\chapter*{Exercises pt. 1} \label{ex 1}
\vspace*{-0.25in}
\begin{exercise}
    {1.1.}Build a cipher wheel as illustrated in Figure 1.1, but with an inner wheel that rotates, and use it to complete the following tasks. 
    \begin{enumerate}
        \item Encrypt the following plaintext using a rotation of 11 clockwise.
        \begin{center}
            ``A page of history is worth a volume of logic.''
        \end{center}
        \item Decrypt the following message, which was encrypted with a rotation of 7 clockwise.
        \begin{center}
            \texttt{AOLYLHYLUVZLJYLAZILAALYAOHUAOLZLJYLALZAOHALCLYFIVKFNBLZZLZ}
        \end{center}
    \end{enumerate}
\end{exercise}

\sol{
\begin{enumerate}
    \item \texttt{L ALRP ZQ STDEZCJ TD HZCES L GZWFXP ZQ WZRTN.}
    \item THERE ARE NO SECRETS BETTER THAN THE SECRETES [sic] THAT EVERY BODY GUESSES \\
    In the encrypted text, ``Secrets'' is \texttt{ZLJYLAZ}. Then, they use an incorrect spelling of the word, \texttt{ZLJYLALZ}, of which has an extra `e' in it. That is what the ``[sic]'' is for. 
\end{enumerate}
}

\begin{exercise}
    {1.2.}Decrypt each of the following Caesar encryptions by trying the various possible
shifts until you obtain readable text.
    \begin{enumerate}
        \item \texttt{LWKLQNWKDWLVKDOOQHYHUVHHDELOOERDUGORYHOBDVDWUHH}
        \item \texttt{UXENRBWXCUXENFQRLQJUCNABFQNWRCJUCNAJCRXWORWMB}
    \end{enumerate}
\end{exercise}

\sol{
\begin{enumerate}
    \item I THINK THAT I SHALL NEVER SEE A BILLBOARD LOVELY AS A TREE
    \item LOVE IS NOT LOVE WHICH ALTERS WHEN IT ALTERATION FINDS
\end{enumerate}
}


\begin{exercise}
    {1.3.}For this exercise, use the simple substitution table given in Table 1.11.
    \begin{enumerate}
        \item Encrypt the plaintext message:
        \center{The gold is hidden in the garden.}
    \end{enumerate}
\end{exercise}

\sol{
\begin{enumerate}
    \item \texttt{IBX FEPA QL BQAAXW QW IBX FSVAXW}
\end{enumerate}
}

\begin{exercise}
    {1.4.} Each of the following messages has been encrypted using a simple substitution cipher. Decrypt them. For your convenience, we have given you a frequency table and a list of the most common bigrams that appear in the ciphertext. (If you do not want to recopy the ciphertexts by hand, they can be downloaded or printed from the web site listed in the preface.)
    \begin{enumerate}
        \item “A Piratical Treasure”
        \begin{flushleft}
            \texttt{JNRZR BNIGI BJRGZ IZLQR OTDNJ GRIHT USDKR ZZWLG OIBTM NRGJN} \\
            \texttt{IJTZJ LZISJ NRSBL QVRSI ORIQT QDEKJ JNRQW GLOFN IJTZX QLFQL} \\
            \texttt{WBIMJ ITQXT HHTBL KUHQL JZKMM LZRNT OBIMI EURLW BLQZJ GKBJT} \\
            \texttt{QDIQS LWJNR OLGRI EZJGK ZRBGS MJLDG IMNZT OIHRK MOSOT QHIJL} \\ 
            \texttt{QBRJN IJJNT ZFIZL WIZTO MURZM RBTRZ ZKBNN LFRVR GIZFL KUHIM} \\ 
            \texttt{MRIGJ LJNRB GKHRT QJRUU RBJLW JNRZI TULGI EZLUK JRUST QZLUK} \\ 
            \texttt{EURFT JNLKJ JNRXR S}
        \end{flushleft}
    \end{enumerate}
\end{exercise}

\sol{
\begin{enumerate}
    \item THESE CHARACTERS AS ONE MIGHT READILY GUESS FORM A CIPHER THAT IS TO SAY THEY CONVEY A MEANING BUT THEN FROM WHAT IS KNOWN OF CAPTAIN KIDD I COULD NOT SUPPOSE HIM CAPABLE OF CONSTRUCTING ANY OF THE MORE ABSTRUSE CRYPTOGRAPHS I MADE UP MY MIND AT ONCE THAT THIS WAS OF A SIMPLE SPECIES SUCH HOW EVER AS WOULD APPEAR TO THE CRUDE INTELLECT OF THE SAILOR ABSOLUTELY INSOLUBLE WITHOUT THE KEY    
\end{enumerate}
\href{https://planetcalc.com/8047/}{Solver}.

}

\begin{exercise}
    {1.5.}Suppose that you have an alphabet of 26 letters. 
    \begin{enumerate}
        \item How many possible simple substitution ciphers are there?
        \item A letter in the alphabet is said to be fixed if the encryption of the letter is the letter itself. \textit{Show an example of how the pieces work together.}
    \end{enumerate}
\end{exercise}

\sol{
\begin{enumerate}
    \item $26!$
    \item This is the formula used for solving for derangements (where $n$ is the number of elements in the set, and $!n$ is the number of derangements [Definition: A permutation with no fixed points]): $!n=n!\sum_{i=0}^{n}{\frac{(-1)^{i}}{i!}}$. From \href{https://en.wikipedia.org/wiki/Derangement}{Wikipedia}. For $n = 2$, we can run through the following:
    \begin{enumerate}[label=(\arabic*)]
        \item $i = 0$: $\frac{(-1)^0}{0!} = 1$
        \item $i = 1$: $\frac{(-1)^1}{1!} = -1$
        \item $i = 2$: $\frac{(-1)^2}{2!} = 0.5$
    \end{enumerate}
    Sum them together: $1 - 1 + 0.5 = 0.5$. Now we can get $!2$: $$!2 = 2! \times (1 - 1 + 0.5) = 2 \times 0.5 = 1$$
\end{enumerate}
}

\begin{exercise}
    {1.6.}Let $a, b, c \in Z$. Use the definition of divisibility to directly prove the following properties of divisibility. (This is Proposition 1.4.)
    \begin{enumerate}
        \item If $a \mid b$ and $b \mid c$, then $a \mid c$.
        \item If $a \mid b$ and $b \mid a$, then $a = \pm b$.
        \item If $a \mid b$ and $a \mid c$, then $a \mid (b + c)$ and $a \mid (b - c)$.
    \end{enumerate}
\end{exercise}

\sol{
    \begin{enumerate}
        \item Let $a,b,c \in \Z$ such that $a \mid b$ and $b \mid c$. We know there exists an $n \in \Z$ such that $a \times n = b$. Similarly, $b \mid c$ means there exists an $k \in \Z$ such that $b\times k = c$. We can use the commutative property to show that:
        \begin{align*}
            k(an) &= (b)k \\
            ank &= bk \\
            bk &= c \\
            a(nk) &= c \\
            a &\mid c
        \end{align*}
        \item From the problem statement, we can intuitively ascertain that because both $a,b$ divide each other, then it must be the case that they are the same number. Moreover, because the criteria for dividing is not pertinent to whether the quotient is negative or positive, this number can be positive or negative. Now that we have an idea of what we are trying to accomplish, we can begin the proof: Let $a,b \in \Z$ such that $a \mid b$ and $b \mid a$. We know there exists an $n \in \Z$ such that $a \times n = b$. Similarly, $b \mid a$ means there exists an $k \in \Z$ such that $b\times k = a$. Then, we can utilize substitution to get the following:
        \begin{align*}
            bk &= a \\
            (an)k &= a 
        \end{align*} From here, we know that because $a$ is present on both sides of the equation, we should divide by $a$ to simplify. Thus, consider the following two cases.
        \begin{itemize}
            \item \textbf{Case 1:} $a \ne 0$ \\
            Since $a$ is not zero, we can divide both sides by $a$ to get $nk = 1$. Since, $n,k \in \Z$, we do not need to worry about fractional reciprocals. Instead, we know from the identity property of multiplication, that $n,k$ must both be $\pm 1$. \begin{itemize}
                \item \textbf{$n,k$ are both $+1$:} Then, $a = b \times 1 \text{ and } b = a \times 1$. Which simplifies to $a = b$ in both cases.
                \item \textbf{$n,k$ are both $-1$:} Then, $a = b \times (-1) \text{ and } b = a \times (-1)$ Which simplifies to $a = -b$ in both cases.
            \end{itemize}
            \item \textbf{Case 2:} $a = 0$ \\
            If $a = 0$, then $b = 0 \times n = 0$ and $0 = b \times k = 0$. Therefore, $a = b = 0$, and $a = \pm b$ is still true.
        \end{itemize}
        We have shown that in either case if $a \mid b$ and $b \mid a$, then $a = \pm b$.
        \item Because $a$ needs to be divisible by both $b$ and $c$, we know that there must exist an $n,k \in \Z$ such that $b = an$ and $c = ak$. Our goal is to get to the form, $a \times \text{ some integer } = b + c$ and $a \times \text{ some integer } = b - c$ so we can use the definition of divides to help us out here. Therefore, let us consider both $b + c$ and $b - c$ in two separate cases:
        \begin{itemize}
            \item \textbf{Case 1:} $b + c$
            \begin{align*}
                b + c &= (an) + (ak) \\
                b + c &= a(n + k) \\
                a &\mid (b + c) 
            \end{align*}
            \item \textbf{Case 2:} $b - c$
            \begin{align*}
                b - c &= (an) - (ak) \\
                b - c &= a(n - k) \\
                a &\mid (b - c) 
            \end{align*}
        \end{itemize}
        We have shown that if $a \mid b$ and $a \mid c$, then $a \mid (b + c)$ and $a \mid (b - c)$.
    \end{enumerate}
}

\begin{exercise}
    {1.7.}Use a calculator and the method described in Remark 1.9 to compute the following quotients and remainders.
    \begin{enumerate}
        \item 34787 divided by 353.
        \item 238792 divided by 7843.
    \end{enumerate}
\end{exercise}

\sol{
    \begin{enumerate}
        \item $a = 34787$ and $b = 353$. Then $a / b \approx 98.54674220$, so $q = 98$ and $r = a - b \cdot q = 34787 - 353 \cdot 98 = 193$.
        \item $a = 238792$ and $b = 7843$. Then $a / b \approx 30.446512$, so $q = 30$ and $r = a - b \cdot q = 238792 - 7843 \cdot 30 = 3502$.
    \end{enumerate}
}

\begin{exercise}
    {1.9.}Use the Euclidean algorithm to compute the following greatest common divisors.
    \begin{enumerate}
        \item $\gcd(291, 252)$.
        \item $\gcd(16261, 85652)$.
    \end{enumerate}
\end{exercise}

\sol{
\begin{enumerate}
    \item $\gcd(291,252)$
    \begin{enumerate}[label=(\arabic*)]
        \item $r_0 = 291$, $r_1 = 252$.
        \item $i = 1$.
        \item Divide $r_0$ by $r_1$ to get a quotient, $q_1$ and a remainder, $r_{2}$:
        \begin{align*}
            291 / 252 &= 1 = q_1 \\
            291 - (252 \times 1) &= 39 = r_2
        \end{align*}
        \item $r_2 \ne 0$. So, we continue.
        \item $i = 2 + 1 = 3$.
    \end{enumerate}
    \pfs
    \begin{enumerate}[label=(\arabic*)]
        \setcounter{enumii}{2}
        \item Divide $r_1$ by $r_2$ to get quotient, $q_2$ and a remainder, $r_3$:
        \begin{align*}
            252 / 39 &= 6 = q_2 \\
            252 - (39 \times 6) &= 18 = r_3
        \end{align*}
        \item $r_3 \ne 0$. So, we continue.
        \item $i = 3 + 1 = 4$.
    \end{enumerate}
    \pfs
    \begin{enumerate}[label=(\arabic*)]
        \setcounter{enumii}{2}
        \item Divide $r_2$ by $r_3$ to get quotient $q_3$ and a remainder, $r_4$:
        \begin{align*}
            39 / 18 &= 2 = q_3 \\
            39 - (18 \times 2) &= 3 = r_4
        \end{align*}
        \item $r_4 \ne 0$. So, we continue.
        \item $i = 3 + 1 = 5$.
    \end{enumerate}
    \pfs
    \begin{enumerate}[label=(\arabic*)]
        \setcounter{enumii}{2}
        \item Divide $r_3$ by $r_4$ to get quotient $q_4$ and a remainder, $r_5$:
        \begin{align*}
            18 / 3 &= 6 = q_4 \\
            18 - (3 \times 6) &= 0 = r_5
        \end{align*}
        \item $r_4 = 0$. So, we stop.
    \end{enumerate}
    We have found that the greatest common divisor is $3$.
    \item \textit{To cut back on paper, I am going to avoid reiterating the steps of 4 and 5. If there is a continuation in the enumeration process, then $r_i \ne 0$, and the process needs to continue}: $\gcd(16261,85652) \Rightarrow \gcd(85652,16261)$.
    \begin{enumerate}[label=(\arabic*)]
    \item
    \begin{align*}
        85652 / 16261 &= 5 = q_1 \\
        85652 - (16261 \times 5) &= 4347 = r_2
    \end{align*}
    \item 
    \begin{align*}
        16261 / 4347 &= 3 = q_2 \\
        16261 - (4347 \times 3) &= 3220 = r_3
    \end{align*}
    \item 
    \begin{align*}
        4347 / 3220 &= 1 = q_3 \\
        4347 - (3220 \times 1) &= 1127 = r_4
    \end{align*}
    \item 
    \begin{align*}
        3220 / 1127 &= 2 = q_4 \\
        3220 - (1127 \times 2) &= 966 = r_5
    \end{align*}
    \item 
    \begin{align*}
        1127 / 966 &= 1 = q_5 \\
        1127 - (966 \times 1) &= 161 = r_6
    \end{align*}
    \item 
    \begin{align*}
        966 / 161 &= 6 = q_6 \\
        966 - (161 \times 6) &= 0 = r_7 
    \end{align*}
    \end{enumerate}
    We have found that $\gcd(85652,16261) = 161$.
    
\end{enumerate}
}

\begin{exercise}
    {1.10.} For each of the $\gcd(a, b)$ values in Exercise 1.9, use the extended Euclidean
algorithm (Theorem 1.11) to find integers $u$ and $v$ such that $au + bv = \gcd(a, b)$.
\end{exercise}

\sol{
\begin{enumerate}
    \item We need to solve for the various $u_i$ and $v_i$. We will start at $i = 2$. \\
    \begin{tabular}{c|c|c|c|c|c|c}
        $i$ & $u_i$ & \textit{Formula} & \textit{Evaluation} & $v_i$ & \textit{Formula} & \textit{Evaluation} \\ \hline
        \textit{2} & $1$ & $u_0 - q_1 \times u_1$ & $1 - 1 \times 0$ & $-1$ & $v_0 - q_1 \times v_1$ & $0 - 1 \times 1$ \\ \hline
        \textit{3} & $-6$ & $u_1 - q_2 \times u_2$ & $0 - 6 \times 1$ & $7$ & $v_1 - q_2 \times v_2$ & $1 - 6 \times (-1)$ \\ \hline
        \textit{4} & $13$ & $u_2 - q_3 \times u_3$ & $1 - 2 \times (-6)$ & $-15$ & $v_2 - q_3 \times v_3$ & $-1 - 2 \times 7$ \\
    \end{tabular} \\

    \pfs
    
    Thus, we can now fill out the table in full:
    \begin{tabular}{c|c|c|c|c|c}
        \textit{i} & \textit{$r_i$} & \textit{$q_i$} & $r_{i + 1}$ & \textit{$u_i$} & \textit{$v_i$} \\ \hline
        \textit{0} & $291$ & $-$ & $-$ & $1$ & $0$ \\ \hline
        \textit{1} & $252$ & $1$ & $39$ & $0$ & $1$ \\ \hline
        \textit{2} & $39$ & $6$ & $18$ & $1$ & $-1$ \\ \hline
        \textit{3} & $18$ & $2$ & $3$ & $-6$ & $7$ \\ \hline
        \textit{4} & $3$ & $6$ & $0$ & $13$ & $-15$ \\
    \end{tabular} \\

    Now, we need to solve: $au + bv = \gcd(a,b) \Rightarrow 291(13) + 252(-15) = 3$. $3$ matches the gcd that we found in Exercise 1.9, so this is the correct solution.
    
    \item 
    
    \begin{tabular}{c|c|c|c|c|c}
        \textit{i} & \textit{$r_i$} & \textit{$q_i$} & $r_{i + 1}$ & \textit{$u_i$} & \textit{$v_i$} \\ \hline
        \textit{0} & $85652$ & $-$ & $-$ & $1$ & $0$ \\ \hline
        \textit{1} & $16261$ & $5$ & $4347$ & $0$ & $1$ \\ \hline
        \textit{2} & $4347$ & $3$ & $3220$ & $1$ & $-5$ \\ \hline
        \textit{3} & $3220$ & $1$ & $1127$ & $-3$ & $16$ \\ \hline
        \textit{4} & $1127$ & $2$ & $966$ & $4$ & $-21$ \\ \hline
        \textit{5} & $966$ & $1$ & $161$ & $-11$ & $58$ \\ \hline
        \textit{6} & $161$ & $6$ & $0$ & $15$ & $-79$ \\
    \end{tabular} \\

    $85652(15) + 16261(-79) = 161$.    
\end{enumerate}
}

\begin{exercise}
    {1.11.} Let $a$ and $b$ be positive integers.
    \begin{enumerate}
        \item Suppose that there are integers $u$ and $v$ satisfying $au + bv = 1$. Prove that $\gcd(a, b) = 1$.
    \end{enumerate}
\end{exercise}

\expf{
\begin{enumerate}
    \item Suppose there are integers $a,b,u,v$ such that $av + bv = 1$. Assume $d \in \Z$ such that $d = \gcd(a,b)$. Since $d$ divides both $a$ and $b$ by definition of common divisor, it must also divide $av$ and $bv$ by definition of divisibility. Moreover, because $au + bv = 1$ and $d$ is a common divisor of both $av$ and $bv$, it must also divide $1$ by Proposition 1.4 (c). Then, the only positive integer that divides $1$ is $1$ itself, so it must be the case that $d = 1$. Therefore, since $d = 1$ and $\gcd(a,b) = d$, it follows that $\gcd(a,b) = 1$.
\end{enumerate}
}

\begin{exercise}
    {1.14.}Let $m \geq 1$ be an integer and suppose that $$a_1 \equiv a_2 \ (\text{mod } m) \text{ and } b_1 \equiv b_2 \ (\text{mod } m).$$ Prove that $$a_1 \pm b_1 \equiv a_2 \pm b_2 \ (\text{mod } m) \text{ and } a_1 \cdot b_1 \equiv a_2 \cdot b_2 \ (\text{mod } m).$$ (This is Proposition 1.13(a).)
\end{exercise}

\expf{
    Let $m \geq 1$ be an integer and suppose that $a_1 \equiv a_2 \ (\text{mod } m) \text{ and } b_1 \equiv b_2 \ (\text{mod } m)$. From the definition of modulo, we know the difference of $a_1 - a_2$ and $b_1 - b_2$ is divisible by $m$. 
    \begin{itemize}
        \item \textbf{Addition:} We want to show that $a_1 + b_1 \equiv a_2 + b_2$. So, our goal is to achieve $m \mid ((a_1 + b_1) - (a_2 + b_2))$. Thus, consider $(a_1 + b_1) - (a_2 + b_2)$. We can distribute the minus sign to get $(a_1 - a_2) + (b_1 - b_2)$. From Proposition 1.4 (c), because we know that $m \mid (a_1 - a_2)$ and $m \mid (b_1 - b_2)$, we can write this as $m \mid ((a_1 - a_2) + (b_1 - b_2))$ which implies $m \mid ((a_1 + b_1) - (a_2 + b_2))$. This shows $a_1 + b_1 \equiv a_2 + b_2 \ (\text{mod } m)$.
        \item \textbf{Subtraction:} Similarly to addition, we want to show $a_1 - b_1 \equiv a_2 - b_2 \ (\text{mod } m)$. Thus, consider $(a_1 - b_1) - (a_2 - b_2) $ which implies $ (a_1 - a_2) - (b_1 - b_2)$. From Proposition 1.4 (c), because we know that $m \mid (a_1 - a_2)$ and $m \mid (b_1 - b_2)$, we can write this as $m \mid ((a_1 - a_2) - (b_1 - b_2))$ which implies $m \mid ((a_1 - b_1) - (a_2 - b_2))$, so $a_1 - b_1 \equiv a_2 - b_2 \ (\text{mod } m)$.
    \end{itemize}
    Therefore we have shown $a_1 \pm b_1 \equiv a_2 \pm b_2 \ (\text{mod } m)$
    \begin{itemize}
        \item \textbf{Product} We want to show that $a_1 \cdot a_2 \equiv a_2 \cdot b_2 \ (\text{mod } m)$. Thus, consider $a_1 \cdot b_1 - a_2 \cdot b_2$:
    \begin{align*}
        a_1 \cdot b_1 - a_2 \cdot b_2 &= a_1 \cdot b_1 - a_1 \cdot b_2 + a_1 \cdot b_2 - a_2 \cdot b_2 \\
        &= a_1 \cdot (b_1 - b_2) + b_2 \cdot (a_1 - a_2)
    \end{align*}
    Because $m \mid (b_1 - b_2)$ and $m \mid (a_1 - a_2)$, we know from the definition of division that when we multiply those numbers by an integer like $a_1$ and $b_2$, $m$ still divides the expression. Hence, $m \mid (a_1 \cdot (b_1 - b_2))$ and $m \mid (b_2 \cdot (a_1 - a_2))$. Therefore, $m \mid (a_1 \cdot b_1 - a_2 \cdot b_2)$ and $a_1 \cdot b_1 \equiv a_2 \cdot b_2 \ (\text{mod } m)$.
    \end{itemize}
}

\begin{exercise}
    {1.16.}Do the following modular computations. In each case, fill in the box with an integer between 0 and $m - 1$, where $m$ is the modulus.
    \begin{enumerate}
        \item $347 + 513 \equiv \boxed{\phantom{000}} \ (\text{mod } 763)$.
        \setcounter{enumi}{2}
        \item $153 \cdot 287 \equiv \boxed{\phantom{000}} \ (\text{mod } 353)$
        \setcounter{enumi}{4}
        \item $5327 \cdot 6135 \cdot 7139 \cdot 2187 \cdot 5219 \cdot 1873 \equiv \boxed{\phantom{000}} \ (\text{mod } 8157)$ (\textit{Hint:} After each multiplication, reduce modulo $8157$ before doing the next multiplication.)
        \setcounter{enumi}{6}
        \item $373^6 \equiv \boxed{\phantom{000}} \ (\text{mod } 581)$.
    \end{enumerate}
\end{exercise}

\sol{
    \begin{enumerate}
        \item $347 + 513 \ (\text{mod } 763) = 97$
        \setcounter{enumi}{2}
        \item $153 \cdot 287 \ (\text{mod } 353) = 139$
        \setcounter{enumi}{4}
        \item 
        \begin{align*}
            5327 \cdot 6135 \ (\text{mod } 8157) &= 4203 \\
            4203 \cdot 7139 \ (\text{mod } 8157) &= 3771 \\
            3771 \cdot 2187 \ (\text{mod } 8157) &= 450 \\
            450 \cdot 5219 \ (\text{mod } 8157) &= 7491 \\
            7491 \cdot 1873 \ (\text{mod } 8157) &= \boxed{603}
        \end{align*}
        \setcounter{enumi}{6}
        \item $373^6 \ (\text{mod } 581) = 463$
    \end{enumerate}
}

\begin{exercise}
    {1.17.}Find all values of $x$ between $0$ and $m - 1$ that are solutions of the following congruences. (\textit{Hint:} If you can’t figure out a clever way to find the solution(s), you can just substitute each value $x = 1$, $x = 2,\dots$, $x = m - 1$ and see which ones work.)
    \begin{enumerate}
        \item $x + 17 \equiv 23 \ (\text{mod } 37)$.
        \setcounter{enumi}{2}
        \item $x^2 \equiv 3 \ (\text{mod } 11)$
        \setcounter{enumi}{6}
        \item $x \equiv 1 \ (\text{mod } 5)$ and also, $x \equiv 2 \ (\text{mod } 7)$. (Find all solutions modulo 35, that is, find the solutions satisfying $0 \leq x \leq 34$.
    \end{enumerate}
\end{exercise}

\sol{
    \begin{enumerate}
        \item $x = 6$
        \setcounter{enumi}{2}
        \item We know that $x$ cannot be any number who's square does not exceed 11 because we cannot square a number to get $3$ (other than $\sqrt{3}$), but these are only the integers, so we cannot use that). Hence, $x \ne 1,2,3$ because we know that the results of these, $1^2 = 1$, $2^2 = 4$, $3^2 = 9$ are not equivalent to 3. Let's try some more values: $x = 4$: $4^2 = 16 \ (\text{mod } 11) = 5 \ \slashed{\equiv} \ 3$; $x = 5$: $5^2 = 25 \ (\text{mod } 11) = 3 \equiv 3$; $x = 6$: $6^2 = 36 \ (\text{mod } 11) = 3 \equiv 3$; $x = 7$: $7^2 = 49 \ (\text{mod } 11) = 5 \ \slashed{\equiv} \ 3$; $x = 8$: $8^2 = 64 \ (\text{mod } 11) = 9 \ \slashed{\equiv} \ 3$; $x = 9$: $9^2 = 81 \ (\text{mod } 11) = 4 \ \slashed{\equiv} \ 3$; $x = 10$: $10^2 = 100 \ (\text{mod } 11) = 1 \ \slashed{\equiv} \ 3$. 

        Thus, we have it that $x = 5,6$.  
        \setcounter{enumi}{6}
        \item Because we have $x \equiv 1 \ (\text{mod } 5)$, verifying correct $x$'s are straightforward. All we need to check for is if a multiple of $5 + 1$ satisfies $x \equiv 2 \ (\text{mod } 7)$. Thus, the only solution is $x = 16$
    \end{enumerate}
}

\begin{exercise}
    {1.19.}Prove that if $a_1$ and $a_2$ are units modulo $m$, then $a_1a_2$ is a unit modulo $m$.
\end{exercise}

\expf{
    Suppose $a_1$ and $a_2$ are units modulo $m$. This means $a \in \Z/m\Z \colon gcd(a,m) = 1$. In other words, $a_1b_1 \equiv 1 \ (\text{mod } m)$ and $a_2b_2 \equiv 1 \ (\text{mod } m)$ for some $b_1, b_2 \in \Z$. When we multiply the equations together, we get $(a_1b_1)(a_2b_2) \equiv 1 \ (\text{mod } m)$ which can be rewritten as $(a_1a_2)(b_1b_2) \equiv 1 \ (\text{mod } m)$. We can multiply $b_1$ and $b_2$ to get an integer $b_3$. Thus, when we multiply $a_1a_2$ by $b_3$ and get $1$, we have shown that $b_3$ is a multiplicative inverse, and $a_1a_2$ is a unit modulo $m$.  
}

\begin{exercise}
    {(Additional).}Decide whether each of the following is a group:
    \begin{enumerate}
        \item \texttt{All 2x2 matrices with real number entries} with operation matrix addition
        \item \texttt{All 2x2 matrices with real number entries} with operation matrix multiplication
    \end{enumerate}

\end{exercise}

\sol{
    \begin{enumerate}
        \item \textbf{Matrix Addition:} \cmark
        \begin{enumerate}[label=(\arabic*)]
            \item \textbf{Closure:} For addition to work between matrices, they must be of dimension $2 \times 2 + 2 \times 2$. Therefore, the dimensions do not change, and it is closed.
            \item \textbf{Associativity:} $2 \times 2$ matrix addition is associative, as it inherits this property from the properties of matrices.
            \item \textbf{Identity Element:} We can add a matrix $Z$ that consists of only 0s to a matrix $A$, and matrix $A$ will remain unchanged.
            \item \textbf{Inverse Element:} True. Consider the matrices, 
            $\begin{bmatrix}
                a & b \\
                c & d
            \end{bmatrix}$, and 
            $\begin{bmatrix}
                -a & -b \\
                -c & -d 
            \end{bmatrix}$. When we add these two together we get
            $\begin{bmatrix}
                0 & 0 \\
                0 & 0
            \end{bmatrix}$. This shows that $2\times 2$ matrices have additive inverses.


            
        \end{enumerate}
        \item \textbf{Matrix Multiplication:} \xmark
        \begin{enumerate}[label=(\arabic*)]
            \item \textbf{Closure:} The dimensions will stay the same during multiplication because it is an $n \times n$ matrix. 
            \item \textbf{Associativity:} $2 \times 2$ matrix multiplication is associative, as it inherits this property from the properties of matrices.
            \item \textbf{Identity Element:} True. Consider the identity matrix, $I_2 =
            \begin{bmatrix}
                1 & 0 \\
                0 & 1
            \end{bmatrix}$. When we multiply a matrix $\begin{bmatrix}
                a & b \\
                c & d
            \end{bmatrix}$ by $I$, we get $$\begin{bmatrix}
                a & b \\
                c & d
            \end{bmatrix} \cdot
            \begin{bmatrix}
                1 & 0 \\
                0 & 1
            \end{bmatrix} = \begin{bmatrix}
                a & b \\
                c & d
            \end{bmatrix},$$
            \item \textbf{Inverse Element:} False. Matrices with a non-zero determinant fail this criteria. Consider the matrix $B =$
            $\begin{bmatrix}
                1 & 2 \\
                3 & 4 
            \end{bmatrix}$. The determinant would be $\det((1)(4) - (2)(3) = -2$. Therefore, this matrix would not have an inverse. 
        \end{enumerate}
    \end{enumerate}
}

\begin{exercise}
    {(Additional).}Is \texttt{All 2x2 matrices with real number entries} a ring with operations matrix addition and matrix multiplication? Justify your answer.
\end{exercise}

\sol{ \cmark
    \begin{enumerate}[label=(\arabic*)]
        \item \textbf{Additive Closure:} True. (See the previous exercise (a), (1)).
        \item \textbf{Additive Associativity:} True. Inherited from the properties of matrices.
        \item \textbf{Additive Identity:} True. (See the previous exercise (a), (3)).
        \item \textbf{Additive Inverse:} True. You can take the difference between a matrix and its inverted duplicate (e.g., $-[A]$) and get $0$.
        \item \textbf{Multiplicative Closure:} True. (See the previous exercise (b), (4)).
        \item \textbf{Distributive Property:} True. While this will pose an error on a calculator, you can do the equivalent: $[A]([B] + [C]) = [A][B] + [A][C]$. This is true because you are still taking the summation of each $a_{ij}$, $b_{ij}$ and $c_{ij}$. 
    \end{enumerate}
}

\chapter*{Exercises pt. 2} \label{ex 2}
\vspace*{-0.25in}
\begin{exercise}
    {1.26}Use the square-and-multiply algorithm described in Section 1.3.2, or the more efficient version in Exercise 1.25, to compute the following powers.
    \begin{enumerate}
        \item \(17^{183} \ (\text{mod } 256).\)
        \item \(2^{477} \ (\text{mod } 1000)\)
    \end{enumerate}
\end{exercise}


\ (\text{mod } 256)


\sol{ \hfill
    \begin{enumerate}
        \item \(a = 17\), \(b = 1\), \(A = 183\), and \(N = 256\). We will now begin the loop because \(A > 0\)
              \begin{itemize}
                  \item \(A = 183\) is odd. Therefore, \(b = b \cdot a \ (\text{mod } 256) = 1 \cdot 17 \ (\text{mod } 256) = 17\); and \(a = a^2 \ (\text{mod } 256) = 17^2 \ (\text{mod 256}) = 33\). \(A = \lfloor 182 / 2 \rfloor = 91\).
                  \item \(A = 91\) is odd. Therefore, \(b = b \cdot a \ (\text{mod } 256) = 17 \cdot 33 \ (\text{mod } 256) = 49\); and \(a = a^2 \ (\text{mod } 256) = 33^2 \ (\text{mod 256}) = 65\). \(A = \lfloor 91 / 2 \rfloor = 45\).
                  \item \(A = 45\) is odd. Therefore, \(b = b \cdot a \ (\text{mod } 256) = 49 \cdot 65 \ (\text{mod } 256) = 113\); and \(a = a^2 \ (\text{mod } 256) = 65^2 \ (\text{mod 256}) = 129\). \(A = \lfloor 45 / 2 \rfloor = 22\).
                  \item \(A = 22\) is even. Therefore, \(a = a^2 \ (\text{mod } 256) = 129^2 \ (\text{mod 256}) = 1\). \(A = \lfloor 22 / 2 \rfloor = 11\).
                  \item \(A = 11\) is odd. Therefore, \(b = b \cdot a \ (\text{mod } 256) = 113 \cdot 1 \ (\text{mod } 256) = 113\); and \(a = a^2 \ (\text{mod } 256) = 1^2 \ (\text{mod 256}) = 1\). \(A = \lfloor 11 / 2 \rfloor = 5\).
                  \item \(A = 5\) is odd. Therefore, \(b = b \cdot a \ (\text{mod } 256) = 113 \cdot 1 \ (\text{mod } 256) = 113\); and \(a = a^2 \ (\text{mod } 256) = 1^2 \ (\text{mod 256}) = 1\). \(A = \lfloor 5 / 2 \rfloor = 2\).
                  \item \(A = 2\) is even. Therefore, \(a = a^2 \ (\text{mod } 256) = 1^2 \ (\text{mod 256}) = 1\). \(A = \lfloor 2 / 2 \rfloor = 1\).
                  \item \(A = 1\) is odd. Therefore, \(b = b \cdot a \ (\text{mod } 256) = 113 \cdot 1 \ (\text{mod } 256) = 113\); and \(a = a^2 \ (\text{mod } 256) = 1^2 \ (\text{mod 256}) = 1\). \(A = \lfloor 1 / 2 \rfloor = 0\).
              \end{itemize}
              Since \(A = 0\), we can report the value of \(b\), which is \(113\). Hence, \(17^183 \ (\text{mod }256) = 113\).
        \item \(a = 2\), \(b = 1\), \(A = 477\), \(N = 1000\). We will now begin the loop because \(A > 0\)
              \begin{itemize}
                  \item \(A = 477\) is odd. Therefore, \(b = b \cdot a \ (\text{mod } 1000) = 1 \cdot 2 \ (\text{mod } 1000) = 2\); and \(a = a^2 \ (\text{mod } 1000) = 2^2 \ (\text{mod } 1000) = 4\). \(A = \lfloor 477 / 2 \rfloor = 238\).
                  \item \(A = 238\) is even. Therefore, \(a = a^2 \ (\text{mod } 1000) = 4^2 \ (\text{mod } 1000) = 16\). \(A = \lfloor 238 / 2 \rfloor = 119\).
                  \item \(A = 119\) is odd. Therefore, \(b = b \cdot a \ (\text{mod } 1000) = 2 \cdot 16 \ (\text{mod } 1000) = 32\); and \(a = a^2 \ (\text{mod } 1000) = 16^2 \ (\text{mod } 1000) = 256\). \(A = \lfloor 119 / 2 \rfloor = 59\).
                  \item \(A = 59\) is odd. Therefore, \(b = b \cdot a \ (\text{mod } 1000) = 32 \cdot 256 \ (\text{mod } 1000) = 192\); and \(a = a^2 \ (\text{mod } 1000) = 256^2 \ (\text{mod } 1000) = 536\). \(A = \lfloor 59 / 2 \rfloor = 29\).
                  \item \(A = 29\) is odd. Therefore, \(b = b \cdot a \ (\text{mod } 1000) = 192 \cdot 536 \ (\text{mod } 1000) = 912\); and \(a = a^2 \ (\text{mod } 1000) = 536^2 \ (\text{mod } 1000) = 296\). \(A = \lfloor 29 / 2 \rfloor = 14\).
                  \item \(A = 14\) is even. Therefore, \(a = a^2 \ (\text{mod } 1000) = 296^2 \ (\text{mod } 1000) = 616\). \(A = \lfloor 14 / 2 \rfloor = 7\).
                  \item \(A = 7\) is odd. Therefore, \(b = b \cdot a \ (\text{mod } 1000) = 912 \cdot 616 \ (\text{mod } 1000) = 792\); and \(a = a^2 \ (\text{mod } 1000) = 616^2 \ (\text{mod } 1000) = 456\). \(A = \lfloor 7 / 2 \rfloor = 3\).
                  \item \(A = 3\) is odd. Therefore, \(b = b \cdot a \ (\text{mod } 1000) = 792 \cdot 456 \ (\text{mod } 1000) = 152\); and \(a = a^2 \ (\text{mod } 1000) = 456^2 \ (\text{mod } 1000) = 936\). \(A = \lfloor 3 / 2 \rfloor = 1\).
                  \item \(A = 1\) is odd. Therefore, \(b = b \cdot a \ (\text{mod } 1000) = 152 \cdot 936 \ (\text{mod } 1000) = 272\); and \(a = a^2 \ (\text{mod } 1000) = 936^2 \ (\text{mod } 1000) = 96\). \(A = \lfloor 1 / 2 \rfloor = 0\).
              \end{itemize}
              Since \(A = 0\), we can report the value of \(b\), which is \(272\). Hence, \(2^477 \ (\text{mod }1000) = 272\).
    \end{enumerate}
}

\begin{exercise}
    {1.41}A \textit{transposition cipher} is a cipher in which the letters of the plaintext remain the same, but their order is rearranged. Here is a simple example in which the message is encrypted in blocks of 25 letters at a time. Take the given 25 letters and arrange them in a 5-by-5 block by writing the message horizontally on the lines. For example, the first 25 letters of the message
    \[\texttt{Now is the time for all good men to come to the aid...} \]
    is written as
    \begin{center}
        \vspace*{-1cm}
        \[
            \begin{array}{ccccc}
                \texttt{N} & \texttt{0} & \texttt{W} & \texttt{I} & \texttt{S} \\
                \texttt{T} & \texttt{H} & \texttt{E} & \texttt{T} & \texttt{I} \\
                \texttt{M} & \texttt{E} & \texttt{F} & \texttt{O} & \texttt{R} \\
                \texttt{A} & \texttt{L} & \texttt{L} & \texttt{G} & \texttt{O} \\
                \texttt{O} & \texttt{D} & \texttt{M} & \texttt{E} & \texttt{N} \\
            \end{array}
        \]
    \end{center}
    Now the ciphertext is formed by reading the letters down the columns, which gives the ciphertext
    \[\texttt{NTMAO OHELD WEFLM ITOGE SIRON.}\]
    \begin{enumerate}
        \item Use this transposition cipher to encrypt the first 25 letters of the message
              \[\texttt{Four score and seven years ago our fathers...}\]
        \item The following message was encrypted using this transposition cipher. Decrypt it.
              \[\texttt{WNOOA HTUFN EHRHE NESUV ICEME.}\]
        \item There are many variations on this type of cipher. We can form the letters into a rectangle instead of a square, and we can use various patterns to place the letters into the rectangle and to read them back out. Try to decrypt the following  ciphertext, in which the letters were placed horizontally into a rectangle of some size and then read off vertically by columns.
              \[\texttt{WHNCE STRHT TEOOH ALBAT DETET SADHE}\]
              \vspace*{-1cm}
              \[\texttt{LEELL QSFMU EEEAT VNLRI ATUDR HTEEA}\]
              (For convenience, we’ve written the ciphertext in 5 letter blocks, but that doesn’t necessarily mean that the rectangle has a side of length 5.)
    \end{enumerate}
\end{exercise}

\sol{ \hfill
    \begin{enumerate}
        \item The first 25 letters of the sentence given is, ``Four score and seven year ago''. Written in block form:
              \begin{center}
                  \vspace*{-1cm}
                  \[
                      \begin{array}{ccccc}
                          \texttt{F} & \texttt{O} & \texttt{U} & \texttt{R} & \texttt{S} \\
                          \texttt{C} & \texttt{O} & \texttt{R} & \texttt{E} & \texttt{A} \\
                          \texttt{N} & \texttt{D} & \texttt{S} & \texttt{E} & \texttt{V} \\
                          \texttt{E} & \texttt{N} & \texttt{Y} & \texttt{E} & \texttt{A} \\
                          \texttt{R} & \texttt{S} & \texttt{A} & \texttt{G} & \texttt{O} \\
                      \end{array}
                  \]
              \end{center}
              Now the ciphertext is formed by reading the letters down the columns, which gives the ciphertext
              \[\texttt{FCNER OODNS URSYA REEEG SAVAO}\]
        \item To create a transposition cipher, we can read the sentence from the starting letter of each word from the sentence, and move inward. For example,
              \begin{enumerate}[label=(\arabic*)]
                  \item \boxed{W}NOOA \boxed{H}TUFN \boxed{E}HRHE \boxed{N}ESUV \boxed{I}CEME \(\Rightarrow\) ``When I...''
                  \item W\boxed{N}OOA H\boxed{T}UFN E\boxed{H}RHE N\boxed{E}SUV I\boxed{C}EME \(\Rightarrow\) ``...N The C...''
                  \item WN\boxed{O}OA HT\boxed{U}FN EH\boxed{R}HE NE\boxed{S}UV IC\boxed{E}ME \(\Rightarrow\) ``...ourse...''
                  \item WNO\boxed{O}A HTU\boxed{F}N EHR\boxed{H}E NES\boxed{U}V ICE\boxed{M}E \(\Rightarrow\) ``...Of Hum...''
                  \item WNOO\boxed{A} HTUF\boxed{N} EHRH\boxed{E} NESU\boxed{V} ICEM\boxed{E} \(\Rightarrow\) ``...an Eve.''
              \end{enumerate}
              Putting it all together, we get, ``\href{https://www.archives.gov/founding-docs/declaration-transcript}{When in the course of human eve...}'' (Declaration of Independence)
        \item This cipher is going to be a little difficult as we have to go through a couple of cases to see which pattern makes the most sense. If we consider the number of letters, 60, we can start at the lowest divisor, 2, and work our way up through the other divisors (this way, we won't start with extremely long ciphers that will each take up a whole page on their lonesome). Thus,
              \begin{itemize}
                  \item Reading by 2s (to get a \(2\times 30\) matrix)
                        \begin{center}
                            \vspace*{-1cm}
                            \[
                                \begin{array}{cccccccccccccccc}
                                    \texttt{W} & \texttt{N} & \texttt{E} & \texttt{T} & \texttt{H} & \texttt{T} & \texttt{O} & \texttt{H} & \texttt{L} & \texttt{A} & \texttt{D} & \texttt{T} & \texttt{A} & \texttt{H} & \texttt{L} & \texttt{E} \\
                                    \texttt{H} & \texttt{C} & \texttt{S} & \texttt{R} & \texttt{T} & \texttt{E} & \texttt{O} & \texttt{A} & \texttt{B} & \texttt{T} & \texttt{E} & \texttt{S} & \texttt{D} & \texttt{E} & \texttt{E} & \texttt{L}
                                \end{array}
                            \]
                            \vspace*{-1cm}
                        \end{center}
                        Clearly, even though we do not have all the letters, there are no words being formed. Let's move on.
                  \item Reading by 3s (to get a \(3\times 20\) matrix)
                        \begin{center}
                            \vspace*{-1cm}
                            \[
                                \begin{array}{cccccccccccccccccccc}
                                    \texttt{W} & \texttt{C} & \texttt{T} & \texttt{T} & \texttt{O} & \texttt{A} & \texttt{A} & \texttt{E} & \texttt{T} & \texttt{D} & \texttt{L} & \texttt{L} & \texttt{S} & \texttt{U} & \texttt{E} & \texttt{V} & \texttt{R} & \texttt{T} & \texttt{R} & \texttt{E} \\

                                    \texttt{H} & \texttt{E} & \texttt{R} & \texttt{T} & \texttt{O} & \texttt{L} & \texttt{T} & \texttt{T} & \texttt{S} & \texttt{H} & \texttt{E} & \texttt{L} & \texttt{F} & \texttt{E} & \texttt{A} & \texttt{N} & \texttt{I} & \texttt{U} & \texttt{H} & \texttt{E} \\

                                    \texttt{N} & \texttt{S} & \texttt{H} & \texttt{E} & \texttt{H} & \texttt{B} & \texttt{D} & \texttt{E} & \texttt{A} & \texttt{E} & \texttt{E} & \texttt{Q} & \texttt{M} & \texttt{E} & \texttt{T} & \texttt{L} & \texttt{A} & \texttt{D} & \texttt{T} & \texttt{A}
                                \end{array}
                            \]
                            \vspace*{-1cm}
                        \end{center}
                        There are no words here.
                  \item \(4 \times 15\) matrix
              \end{itemize}
              \begin{center}
                  \vspace*{-1cm}
                  \[
                      \begin{array}{ccccccccccccccc}
                          \texttt{W} & \texttt{E} & \texttt{H} & \texttt{O} & \texttt{L} & \texttt{D} & \texttt{T} & \texttt{H} & \texttt{E} & \texttt{S} & \texttt{E} & \texttt{T} & \texttt{R} & \texttt{U} & \texttt{T} \\

                          \texttt{H} & \texttt{S} & \texttt{T} & \texttt{O} & \texttt{B} & \texttt{E} & \texttt{S} & \texttt{E} & \texttt{L} & \texttt{F} & \texttt{E} & \texttt{V} & \texttt{I} & \texttt{D} & \texttt{E} \\

                          \texttt{N} & \texttt{T} & \texttt{T} & \texttt{H} & \texttt{A} & \texttt{T} & \texttt{A} & \texttt{L} & \texttt{L} & \texttt{M} & \texttt{E} & \texttt{N} & \texttt{A} & \texttt{R} & \texttt{E} \\

                          \texttt{C} & \texttt{R} & \texttt{E} & \texttt{A} & \texttt{T} & \texttt{E} & \texttt{D} & \texttt{E} & \texttt{Q} & \texttt{U} & \texttt{A} & \texttt{L} & \texttt{T} & \texttt{H} & \texttt{A}
                      \end{array}
                  \]
                  \vspace*{-1cm}
              \end{center}
              There we finally have our transcription! ``We hold these truths to be self evident that all men are created equal tha...'' (Declaration of Independence again).
    \end{enumerate}
}

\begin{exercise}
    {Additional Problem}Write down the steps for an algorithm to encrypt a plaintext using a transposition/Scytale cipher using \(n\) columns. \textit{Hint: This should involve some modular arithmetic.}
\end{exercise}

\sol{
    \begin{enumerate}[label=(\arabic*)]
        \item Remove all spaces from the plaintext message.
        \item Let \(i\) be the length of the message after removing spaces.
        \item Find the number of rows \(m\) required for the transposition cipher matrix:
              \[
                  m = \left\lceil \frac{i}{n} \right\rceil
              \]
        \item Create an \(m \times n\) matrix to hold the characters.
        \item For each character in the message, place it into the matrix. Let the index of the current character in the message be \(k\), where \(0 \leq k < i\). Calculate the row and column for this character as follows:
              \[
                  \text{row} = \left\lfloor \frac{k}{n} \right\rfloor
              \]
              \[
                  \text{column} = k \ (\text{mod } n)
              \]
        \item Once all characters are placed into the matrix, read the matrix column by column to form the ciphertext. For each column \(j\), loop through the rows and append each character to the ciphertext in the following order:
              \[
                  \text{character index} = j + n \times \text{row}
              \]
        \item If there are empty cells (i.e., \(n\) does not divide the message length \(i\)), fill these cells with random characters or leave them blank.
        \item Concatenate the characters column by column to form the final ciphertext.
    \end{enumerate}

}

\begin{exercise}
    {1.32}For each of the following primes \( p \) and numbers \( a \), compute \( a^{-1} \mod p \) in two ways: (i) the extended Euclidean algorithm. (ii) Use the fast power algorithm and Fermat’s little theorem. (See Example 1.27.)
    \begin{enumerate}
        \item \( p = 47 \) and \( a = 11 \).
        \item \( p = 587 \) and \( a = 345 \).
    \end{enumerate}
\end{exercise}

\sol{ \hfill
    \begin{enumerate}
        \item \begin{enumerate}[label=(\roman*)]
                  \item For the extended Euclidean algorithm, we can start by filling out this table:
                        \begin{center}
                            \begin{tabular}{c|c|c|c|c|c}
                                \textit{i} & \textit{\(r_i\)} & \textit{\(q_i\)} & \(r_{i + 1}\) & \textit{\(u_i\)}        & \textit{\(v_i\)}          \\ \hline
                                \textit{0} & \(47\)           & \(-\)            & \(-\)         & \(1\)                   & \(0\)                     \\ \hline
                                \textit{1} & \(11\)           & \(4\)            & \(3\)         & \(0\)                   & \(1\)                     \\ \hline
                                \textit{2} & \(3\)            & \(3\)            & \(2\)         & \(1\)                   & \(0 - 4 \cdot 1 = -4\)    \\ \hline
                                \textit{3} & \(2\)            & \(1\)            & \(1\)         & \(0 - 3 \cdot 1 = -4\)  & \(1 - 3 \cdot -4 = 13\)   \\ \hline
                                \textit{4} & \(1\)            & \(2\)            & \(0\)         & \(1 - 1 \times -4 = 4\) & \(-4 - 1 \cdot 13 = -17\)
                            \end{tabular} \\
                        \end{center}
                        Thus, to find  \(a^{-1} \ (\text{mod } p)\), we need an \(x\) such that \(11x \equiv 1 \ (\text{mod } 47)\). From our table, we know that number to be \(-17\) because \(1 = 4 \times 47 - 17 \times 11\). Then, as a positive number \(\text{mod } 47\), we get \(x \equiv -17 \equiv 30 \ (\text{mod } 47)\).
                  \item For Fermat's Little Theorem:
                        \begin{align*}
                            a^{p - 1} & \equiv 1 \ (\text{mod } p)      \\
                            a^{p - 2} & \equiv a^{-1} \ (\text{mod } p) \\
                            11^{45}   & \ (\text{mod } 47)
                        \end{align*}
                        By the fast power algorithm, we need to find the binary representation of 45. Thus, \(45 \ (\text{mod } 2) = 1\), \(22 \ (\text{mod } 2) = 0\), \(11 \ (\text{mod } 2) = 1\), \(5 \ (\text{mod } 2) = 1\), \(2 \ (\text{mod } 2) = 0\), \(1 \ (\text{mod } 2) = 1\). From this, we have the binary representation of \(45_{10}\) as \(101101_2\). Hence, we will need to calculate \(2^5 \cdot 2^3 \cdot 2^2 \cdot 2^1 \cdot 2^0 \Rightarrow 11^{2^5} \cdot 11^{2^3} \cdot 11^{2^0} \Rightarrow 11^{32} \cdot 11^8 \cdot 11^4 \cdot 11 \equiv 17 \cdot 14 \cdot 25 \cdot 11\). Now, lets find the values of these numbers:
                        \begin{align*}
                            11^1    & \equiv 11 \ (\text{mod } 47) \\
                            11^2    & \equiv 27 \ (\text{mod } 47) \\
                            11^4    & \equiv 25 \ (\text{mod } 47) \\
                            11^8    & \equiv 14 \ (\text{mod } 47) \\
                            11^{16} & \equiv 8 \ (\text{mod } 47)  \\
                            11^{32} & \equiv 17 \ (\text{mod } 47)
                        \end{align*}
                        Now, we can calculate \(11^{45} = 11^{32} \cdot 11^8 \cdot 11^4 \cdot 11 \equiv 17 \cdot 14 \cdot 25 \cdot 11 \ (\text{mod } 47)\). Further multiplying we get \(17 \cdot 14 \ (\text{mod } 47) \equiv 3\). Then, \(3 \cdot 25 \ (\text{mod } 47) \equiv 28\). And finally, \(28 \cdot 11 \ (\text{mod } 47) \equiv 30\). Thus confirming our previous answer in (i).
              \end{enumerate}
              %
              %
              %
              %
              %   
        \item \begin{enumerate}[label=(\roman*)]
                  \item For the extended Euclidean algorithm:
                        \begin{center}
                            \begin{tabular}{c|c|c|c|c|c}
                                \textit{i} & \textit{\(r_i\)} & \textit{\(q_i\)} & \(r_{i + 1}\) & \textit{\(u_i\)}          & \textit{\(v_i\)}           \\ \hline
                                \textit{0} & \(587\)          & \(-\)            & \(-\)         & \(1\)                     & \(0\)                      \\ \hline
                                \textit{1} & \(345\)          & \(1\)            & \(242\)       & \(0\)                     & \(1\)                      \\ \hline
                                \textit{2} & \(242\)          & \(1\)            & \(103\)       & \(1\)                     & \(0 - 1 \cdot 1 = -1\)     \\ \hline
                                \textit{3} & \(103\)          & \(2\)            & \(36\)        & \(0 - 1 \cdot 1 = -1\)    & \(1 - 1 \cdot -1 = 2\)     \\ \hline
                                \textit{4} & \(36\)           & \(2\)            & \(31\)        & \(1 - 2 \cdot -1 = 3\)    & \(-1 - 2 \cdot 2 = -5\)    \\ \hline
                                \textit{5} & \(31\)           & \(1\)            & \(5\)         & \(-1 - 2 \cdot 3 = -7\)   & \(2 - 2 \cdot -5 = 12\)    \\ \hline
                                \textit{6} & \(5\)            & \(6\)            & \(1\)         & \(3 - 1 \cdot -7 = 10\)   & \(-5 - 1 \cdot 12 = -17\)  \\ \hline
                                \textit{7} & \(1\)            & \(5\)            & \(0\)         & \(10 - 6 \cdot 10 = -67\) & \(12 - 6 \cdot -17 = 114\)
                            \end{tabular}
                        \end{center}
                        Thus, the inverse of \(345 \ (\text{mod } 587)\) is 114.

                  \item The binary expression for \(585_{10}\) is \(1001001001_2\) (I just used my calculator this time and kept dividing ans by 2 while keeping track of the odd vs even quotients). This leaves us with \(345^{2^9} \cdot 345^{2^6} \cdot 345^{2^3} \cdot 345^{2^0}\).
                        \begin{align*}
                            345^1     & \equiv 345 \ (\text{mod } 587) \\
                            345^2     & \equiv 451 \ (\text{mod } 587) \\
                            345^4     & \equiv 299 \ (\text{mod } 587) \\
                            345^8     & \equiv 177 \ (\text{mod } 587) \\
                            345^{16}  & \equiv 218 \ (\text{mod } 587) \\
                            345^{64}  & \equiv 529 \ (\text{mod } 587) \\
                            345^{128} & \equiv 429 \ (\text{mod } 587) \\
                            345^{256} & \equiv 310 \ (\text{mod } 587) \\
                            345^{512} & \equiv 419 \ (\text{mod } 587) \\
                        \end{align*}
                        Now we can multiply and solve: \(345^{2^9} \cdot 345^{2^6} \cdot 345^{2^3} \cdot 345^{2^0} \Rightarrow 419 \cdot 529 \cdot 177 \cdot 345 \ (\text{mod } 587) = 114\). Therefore, the inverse of \(345 \ (\text{mod } 587)\) is \(114\).
              \end{enumerate}
    \end{enumerate}
}

\begin{exercise}
    {1.34}Recall that \( g \) is called a primitive root modulo \( p \) if the powers of \( g \) give all nonzero elements of \( \mathbb{F}_p \).
    \begin{enumerate}[label=(\alph*)]
        \item For which of the following primes is 2 a primitive root modulo \( p \)?
              \[
                  \begin{array}{llllllll}
                      \text{(i)} & p = 7 & \text{(ii)} & p = 13 & \text{(iii)} & p = 19 & \text{(iv)} & p = 23
                  \end{array}
              \]
        \item For which of the following primes is 3 a primitive root modulo \( p \)?
              \[
                  \begin{array}{llllllll}
                      \text{(i)} & p = 5 & \text{(ii)} & p = 7 & \text{(iii)} & p = 11 & \text{(iv)} & p = 17
                  \end{array}
              \]
        \item Find a primitive root for each of the following primes.
              \[
                  \begin{array}{llllllll}
                      \text{(i)} & p = 23 & \text{(ii)} & p = 29 & \text{(iii)} & p = 41 & \text{(iv)} & p = 43
                  \end{array}
              \]
        \item Find all primitive roots modulo 11. Verify that there are exactly \(\phi(10)\) of them, as asserted in Remark 1.32.
    \end{enumerate}
\end{exercise}

\sol{ \hfill
    \begin{enumerate}
        \item \begin{enumerate}[label=(\roman*)]
                  \item \(p = 7\): \(2^1 \equiv 2\), \(2^2 \equiv 4\), \(2^3 \equiv 1 \ (\text{mod } 7)\). Because \(2\) does not cover every nonzero elements of \( \mathbb{F}_p \), \(2\) is not a primitive root.
                  \item \(p = 13\): \(2^1 \equiv 2\), \(2^2 \equiv 4\), \(2^3 \equiv 8\), \(2^4 \equiv 3\), \(2^5 \equiv 6\), \(2^6 \equiv 12\), \(2^7 \equiv 11\), \(2^8 \equiv 9\), \(2^9 \equiv 5\), \(2^{10} \equiv 10\), \(2^{11} \equiv 7\), \(2^{12} \equiv 1\), \(2\) is a primitive root.
                  \item \(p = 19\): \(2^1 \equiv 2\), \(2^2 \equiv 4\), \(2^3 \equiv 8\), \(2^{4} \equiv 16\), \(2^{5} \equiv 13\), \(2^{6} \equiv 7\), \(2^{7} \equiv 14\), \(2^{8} \equiv 9\), \(2^{9} \equiv 18\), \(2^{10} \equiv 17\), \(2^{11} \equiv 15\), \(2^{12} \equiv 11\), \(2^{13} \equiv 3\), \(2^{14} \equiv 6\), \(2^{15} \equiv 12\), \(2^{16} \equiv 5\), \(2^{17} \equiv 10\), \(2^{18} \equiv 1\), 2 is a primitive root.
                  \item \(p = 23\): \(2^1 \equiv 2\), \(2^2 \equiv 4\), \(2^3 \equiv 8\), \(2^{4} \equiv 16\), \(2^{5} \equiv 9\), \(2^{6} \equiv 18\), \(2^{7} \equiv 13\), \(2^{8} \equiv 3\), \(2^{9} \equiv 6\), \(2^{10} \equiv 12\), \(2^{11} \equiv 1\), \(2\) is not a primitive root.
              \end{enumerate}
        \item \begin{enumerate}[label=(\roman*)]
                  \item \( p = 5 \): \( 3^1 \equiv 3, 3^2 \equiv 4, 3^3 \equiv 2, 3^4 \equiv 1 \), 3 is a primitive root.
                  \item \( p = 7 \): \( 3^1 \equiv 3, 3^2 \equiv 2, \ldots, 3^6 \equiv 1 \), 3 is a primitive root.
                  \item \( p = 11 \): \( 3^1 \equiv 3, 3^2 \equiv 9, 3^3 \equiv 5, \ldots, 3^5 \equiv 1 \), 3 is not a primitive root.
                  \item \( p = 17 \): \( 3^1 \equiv 3, 3^2 \equiv 9, \ldots, 3^{16} \equiv 1 \), 3 is a primitive root.
              \end{enumerate}
        \item \begin{enumerate}[label=(\roman*)]
                  \item \(p = 23\): \(5^1 \equiv 5\),
                        \(5^{2} \equiv 2
                        \), \(5^{3} \equiv 10
                        \), \(5^{4} \equiv 4
                        \), \(5^{5} \equiv 20
                        \), \(5^{6} \equiv 8
                        \), \(5^{7} \equiv 17
                        \), \(5^{8} \equiv 16
                        \), \(5^{9} \equiv 11
                        \), \(5^{10} \equiv 9
                        \), \(5^{11} \equiv 22
                        \), \(5^{12} \equiv 18
                        \), \(5^{13} \equiv 21
                        \), \(5^{14} \equiv 13
                        \), \(5^{15} \equiv 19
                        \), \(5^{16} \equiv 3
                        \), \(5^{17} \equiv 15
                        \), \(5^{18} \equiv 6
                        \), \(5^{19} \equiv 7
                        \), \(5^{20} \equiv 12
                        \), \(5^{21} \equiv 14\),
                        \(5^{22} \equiv 1\), \\
                        \(5\) is a primitive root.
                  \item \(2^1 \equiv 2\),
                        \(2^{2} \equiv 4\),
                        \(2^{3} \equiv 8\),
                        \(2^{4} \equiv 16\),
                        \(2^{5} \equiv 3\),
                        \(2^{6} \equiv 6\),
                        \(2^{7} \equiv 12\),
                        \(2^{8} \equiv 24\),
                        \(2^{9} \equiv 19\),
                        \(2^{10} \equiv 9\),
                        \(2^{11} \equiv 18\),
                        \(2^{12} \equiv 7\),
                        \(2^{13} \equiv 14\),
                        \(2^{14} \equiv 28\),
                        \(2^{15} \equiv 27\),
                        \(2^{16} \equiv 25\),
                        \(2^{17} \equiv 21\),
                        \(2^{18} \equiv 13\),
                        \(2^{19} \equiv 26\),
                        \(2^{20} \equiv 23\),
                        \(2^{21} \equiv 17\),
                        \(2^{22} \equiv 5\),
                        \(2^{23} \equiv 10\),
                        \(2^{24} \equiv 20\),
                        \(2^{25} \equiv 11\),
                        \(2^{26} \equiv 22\),
                        \(2^{27} \equiv 15\),
                        \(2^{28} \equiv 1\), \\
                        \(2\) is a primitive root.
                  \item \(6^{1} \equiv 6\)
                        \(6^{2} \equiv 36\),
                        \(6^{3} \equiv 11\),
                        \(6^{4} \equiv 25\),
                        \(6^{5} \equiv 27\),
                        \(6^{6} \equiv 39\),
                        \(6^{7} \equiv 29\),
                        \(6^{8} \equiv 10\),
                        \(6^{9} \equiv 19\),
                        \(6^{10} \equiv 32\),
                        \(6^{11} \equiv 28\),
                        \(6^{12} \equiv 4\),
                        \(6^{13} \equiv 24\),
                        \(6^{14} \equiv 21\),
                        \(6^{15} \equiv 3\),
                        \(6^{16} \equiv 18\),
                        \(6^{17} \equiv 26\),
                        \(6^{18} \equiv 33\),
                        \(6^{19} \equiv 34\),
                        \(6^{20} \equiv 40\),
                        \(6^{21} \equiv 35\),
                        \(6^{22} \equiv 5\),
                        \(6^{23} \equiv 30\),
                        \(6^{24} \equiv 16\),
                        \(6^{25} \equiv 14\),
                        \(6^{26} \equiv 2\),
                        \(6^{27} \equiv 12\),
                        \(6^{28} \equiv 31\),
                        \(6^{29} \equiv 22\),
                        \(6^{30} \equiv 9\),
                        \(6^{31} \equiv 13\),
                        \(6^{32} \equiv 37\),
                        \(6^{33} \equiv 17\),
                        \(6^{34} \equiv 20\),
                        \(6^{35} \equiv 38\),
                        \(6^{36} \equiv 23\),
                        \(6^{37} \equiv 15\),
                        \(6^{38} \equiv 8\),
                        \(6^{39} \equiv 7\),
                        \(6^{40} \equiv 1\),\\
                        \(6\) is a primitive root.
                  \item \(3^{1} \equiv 3\),
                        \(3^{2} \equiv 9\),
                        \(3^{3} \equiv 27\),
                        \(3^{4} \equiv 38\),
                        \(3^{5} \equiv 28\),
                        \(3^{6} \equiv 41\),
                        \(3^{7} \equiv 37\),
                        \(3^{8} \equiv 25\),
                        \(3^{9} \equiv 32\),
                        \(3^{10} \equiv 10\),
                        \(3^{11} \equiv 30\),
                        \(3^{12} \equiv 4\),
                        \(3^{13} \equiv 12\),
                        \(3^{14} \equiv 36\),
                        \(3^{15} \equiv 22\),
                        \(3^{16} \equiv 23\),
                        \(3^{17} \equiv 26\),
                        \(3^{18} \equiv 35\),
                        \(3^{19} \equiv 19\),
                        \(3^{20} \equiv 14\),
                        \(3^{21} \equiv 42\),
                        \(3^{22} \equiv 40\),
                        \(3^{23} \equiv 34\),
                        \(3^{24} \equiv 16\),
                        \(3^{25} \equiv 5\),
                        \(3^{26} \equiv 15\),
                        \(3^{27} \equiv 2\),
                        \(3^{28} \equiv 6\),
                        \(3^{29} \equiv 18\),
                        \(3^{30} \equiv 11\),
                        \(3^{31} \equiv 33\),
                        \(3^{32} \equiv 13\),
                        \(3^{33} \equiv 39\),
                        \(3^{34} \equiv 31\),
                        \(3^{35} \equiv 7\),
                        \(3^{36} \equiv 21\),
                        \(3^{37} \equiv 20\),
                        \(3^{38} \equiv 17\),
                        \(3^{39} \equiv 8\),
                        \(3^{40} \equiv 24\),
                        \(3^{41} \equiv 29\),
                        \(3^{42} \equiv 1\), \\
                        \(3\) is a primitive root.
              \end{enumerate}
        \item The primitive roots modulo 11 are 2, 6, 7, 8. To check: \\
              \(2^1 \equiv 11 = 2\),
              \(2^2 \equiv 11 = 4\),
              \(2^3 \equiv 11 = 8\),
              \(2^4 \equiv 11 = 5\),
              \(2^5 \equiv 11 = 10\),
              \(2^6 \equiv 11 = 9\),
              \(2^7 \equiv 11 = 7\),
              \(2^8 \equiv 11 = 3\),
              \(2^9 \equiv 11 = 6\),
              \(2^{10} \equiv 11 = 1\),

              2 is a primitive root. \\

              \(6^1 \equiv 11 = 6\),
              \(6^2 \equiv 11 = 3\),
              \(6^3 \equiv 11 = 7\),
              \(6^4 \equiv 11 = 9\),
              \(6^5 \equiv 11 = 10\),
              \(6^6 \equiv 11 = 5\),
              \(6^7 \equiv 11 = 8\),
              \(6^8 \equiv 11 = 4\),
              \(6^9 \equiv 11 = 2\),
              \(6^{10} \equiv 11 = 1\),

              6 is a primitive root. \\

              \(7^1 \equiv 11 = 7\),
              \(7^2 \equiv 11 = 5\),
              \(7^3 \equiv 11 = 2\),
              \(7^4 \equiv 11 = 3\),
              \(7^5 \equiv 11 = 10\),
              \(7^6 \equiv 11 = 4\),
              \(7^7 \equiv 11 = 6\),
              \(7^8 \equiv 11 = 9\),
              \(7^9 \equiv 11 = 8\),
              \(7^{10} \equiv 11 = 1\),

              7 is a primitive root. \\

              \(8^1 \equiv 11 = 8\),
              \(8^2 \equiv 11 = 9\),
              \(8^3 \equiv 11 = 6\),
              \(8^4 \equiv 11 = 4\),
              \(8^5 \equiv 11 = 10\),
              \(8^6 \equiv 11 = 3\),
              \(8^7 \equiv 11 = 2\),
              \(8^8 \equiv 11 = 5\),
              \(8^9 \equiv 11 = 7\),
              \(8^{10} \equiv 11 = 1\),

              8 is a primitive root.
    \end{enumerate}
}



\chapter{The Real Numbers} \label{Real nums}
\vspace*{-0.25in}
\section{Types of Numbers}

The \gls{natural numbers}, $\N$:
\begin{itemize}
    \item No additive inverse.
    \item You can:
    \begin{itemize}
        \item Add,
        \item Multiply
    \end{itemize}
\end{itemize}

The \gls{integers}, $\Z$ are known as a Group (more specifically, a ``ring'').

\begin{itemize}
    \item You can:
    \begin{itemize}
        \item Add,
        \item Multiply,
        \item Subtract
    \end{itemize}
\end{itemize}

The \gls{rational numbers}, $\Q$ are known as a ``Field.''

\begin{itemize}
    \item You can:
    \begin{itemize}
        \item Add,
        \item Subtract,
        \item Multiply,
        \item Divide
    \end{itemize}
\end{itemize}

A problem that rational numbers could not explain: The $45, 45, 90$ triangle had a hypotenuse of $\sqrt{2}$. This did not exist at the time, so it was simply $c^2 = 2$. Therefore, new numbers needed to be invented.

\begin{theorem}
    There does not exist a rational number $r$ such that $r^2 = 2$.
\end{theorem}

\tpf{
    Suppose there exists a rational number $r$ such that $r^2 = 2$. Since $r$ is rational, there exists $p,q \in \Z$ such that $r = \frac{p}{q}$. We can assume the $p$ and $q$ have no common factors. (If not, we can factor out the common factor.) By our assumption, 
    \begin{align*}
        r^2 &= 2 \\
        \frac{p^2}{q^2} &= 2 
    \end{align*}
    It follows that, \begin{align*}
        p^2 = 2q^2
    \end{align*}
    Such that $p^2$ is an even number because if $p$ were odd, then $p^2$ would be odd. There exists $x \in \Z$ such that $p = 2x$. Recall that $p^2 = 2q^2$. Thus\begin{align*}
        (2x)^2 &= 2q^2 \\
        4x^2 &= 2q^2 \\
        2x^2 &= q
    \end{align*} Thus, $q^2$ is even. Hence $q$ is also even. So $p$ and $q$ are both divisible by 2. This contradicts that $p$ and $q$ have no common factors. Thus, our supposition is false. Therefore, there does not exist a rational number $r$ such that $r^2 = 2$ 
}

So we are going to work with a larger set called the real numbers, $\R$. 
\begin{itemize}
    \item $\N \subset \Z \subset \Q \subset \R$
    \item You can:
    \begin{itemize}
        \item Add,
        \item Subtract,
        \item Multiply,
        \item Divide
    \end{itemize}
    \item In other words, all field axioms apply.
    \item Totally ordered set for any $x,y \in \R$. Thus, one of these are true:
    \begin{enumerate}
        \item $x < y$,
        \item $x > y$,
        \item $x = y$
    \end{enumerate}
    \item  Think of it as a number line.
    \item $\Q$ is dense: \\
    If $a,b \in \Q$ with $a \ne b$, there exists $c \in \Q$ which is between $a$ and $b$ such that $a < c < b$. One example is $\frac{a + b}{2}$.
    \item $\Q$ is not \textit{complete}, but $\R$ is.
    \begin{itemize}
        \item \textit{Complete:} Think, ``no gaps.''
    \end{itemize}
\end{itemize}

\section{Preliminaries}

Things to remember from Intro and Discrete.

\begin{table}[htbp]
\begin{tabular}{c||c}
Set Notation & Complement \\
\hline 
$x \in A$ & $A^c$ (not $\overline{A}$) \\
$A \cup B$ & $\R \setminus A$ \\
$A \cap B$
\end{tabular}
\end{table}

\begin{itemize}
\item $\displaystyle \bigcup^\infty_{n=1} A_n = A_1 \cup A_2 \cup A_3 \cup \dots$. \\

\item $\displaystyle \bigcap^\infty_{n=1} A_n = A_1 \cap A_2 \cap \dots$. \\

\item \textbf{\Gls{De Morgan's Laws}} \\
\end{itemize}


\pfs

\subsection{Infinite Unions and Intersections} \hfill

For each $n \in \N$, define $A_n = \{n, n+1, n+2,\dots\} = \{k \in \N \ | \ k \geq n\}$. In other words, each subsequent element in the subset will start at $n$. For example, $A_1 = \{1,2,\dots\}$, whereas $A_5 = \{5,6,\dots\}.$   \\

$\bigcup^\infty_{n=1} A_n = \N$. To show a number $\in \N$ belongs in the set $A_n$, we can start with that, $k \in \N$. Then $k \in A_k$. Thus, $k \in A_k \subseteq \bigcup^\infty_{n=1}A_n$. therefore, $\N \subseteq \bigcup^\infty_{n=1}A_n$. \\

$\bigcap^\infty_{n=1}A_n = \emptyset$. Obviously, we know that the empty set is a subset of $A_n$, but to prove that $\bigcap^\infty_{n=1}A_n$ is a subset of the empty set, we should suppose a $k \in \N$ such that $k \in \bigcap^\infty_{n=1}A_n$. Notice that $k \notin \bigcap^\infty_{n=1}A_n$. So, $\bigcap^\infty_{n=1}A_n = \emptyset$.

\pfs

\subsection{Functions and Notation} \hfill

$f \colon A \rightarrow B$ where $f$ is a function, $A$ is a domain, and $B$ is the co-domain. Thus, $f(x) = y$ such that $x \in A$ and $y \in B$. \\

Some definitions to keep in mind: \\

\textbf{\Gls{The Dirichlet Function}} \\

\textit{[Refer to notepaper for these following definitions]} \\

\textbf{\Gls{image}} \\

Example: $g \ \colon \ \R \rightarrow \R$, when we say $y \in g(A)$ implies $\exists x$ such that $g(x) = y$

\textbf{\Gls{pre-image}} \\

\textbf{\Gls{triangle inequality}}: \\
The most common application: For any $a,b,c \in \R$, $|a - b| \leq |a-c| + |c-b|$, with the intermediate step of $a - b = (a-c) + (c-b)$.

\subsection{Common Strategies for Analysis Proofs} \hfill

\setcounter{BoxCounter}{5}
\begin{theorem}
    Let $a,b \in \R$. Then, $$a = b \text{ if and only if for all } \epsilon > 0, |a-b| < \epsilon.$$
\end{theorem}

\iffpf
    {Assume $a = b$. Let $\epsilon > 0$. Then $|a-b| = 0 < \epsilon$}
    {Assume for all $\epsilon > 0$, $|a-b| < \epsilon$. Suppose $a \ne b$. Then $a - b \ne 0$. So, $|a-b| \ne 0$. Now, Consider $\epsilon_0 = |a-b|$. By our assumption we know that $|a-b| < \epsilon_0$. It is not true that $|a-b| < |a - b|$. Therefore, it must be the case that $a = b$.}
    {Therefore, by showing both sides of the implication accomplish the same thing as the other side, we know that $a = b \text{ if and only if for all } \epsilon > 0, |a-b| < \epsilon.$}
    {xblue}

\subsection{Mathematical Induction} \hfill

\textit{Inductive Hypothesis:} Let $x_1 = 1$. For all $n \in \N$, let $x_{n+1} = \frac{1}{2}x_n + 1$. \\

\textit{Inductive Step:} $x_1 = 1, x_2 = 1.5, x_3 = 1.75, x_4 = 1.875$. \\

\begin{example}
    {Induction} The sequence $(x_n)$ is increasing. In other words, for all $n\in \N$, $x_n \leq x_{n+1}$. 
\end{example}

\epf{
    Suppose the sequence $(x_n)$ is increasing. We will prove this point by using induction. \\
    
    \textbf{Base Case:} We see that $x_1 = 1$ and $x_2 = 1.5$. Thus, $x_1 \leq x_2$. \\
    
    \textbf{Inductive Hypothesis:} For $n \in \N$, assume $x_n \leq x_{n+1}$. \\
    
    \indent \textit{Scratch work:} We want: $x_{n+1} \leq x_{n+2}$. We know: $x_{n+1} = \frac{1}{2}x_{n+1}+1$. \\
    
    \textbf{Inductive Step:} Then $\frac{1}{2}x_n \leq \frac{1}{2}x_{n+1}$. Hence, $\frac{1}{2}x_{n} + 1 \leq \frac{1}{2}x_{n+1} + 1$. Therefore we have proven through induction that, $x_{n+1} \leq x_{n+2}$.
}

\subsection*{Exercises}

\begin{exercise}
    {1.2.3}Decide which of the following represent true statements about
the nature of sets. For any that are false, provide a specific example where the
statement in question does not hold.
\begin{enumerate}
    \item If $A_1 \supseteq A_2 \supseteq A_3 \supseteq A_4\dots$ are all sets containing an infinite number of
elements, then the intersection $\bigcap^\infty_{n=1} A_n$ is infinite as well.
    \item If $A_1 \supseteq A_2 \supseteq A_3 \supseteq A_4\dots$ are all finite, nonempty sets of real numbers,
then the intersection $\bigcap^\infty_{n=1} A_n$ is finite and nonempty.
    \item $A \cap (B \cup C) = (A \cap B) \cup C$
    \item $A \cap (B \cap C) = (A \cap B) \cap C$
    \item $A \cap (B \cup C) = (A \cap B) \cup (A \cap C)$
\end{enumerate}
\end{exercise}

\sol{\hfill
    \begin{enumerate}
        \item This is false. Consider the following as a counterexample: If we define $A_1$ as $A_n = \{n, n+1, n+2,\dots\} = \{k \in \N \ | \ k \geq n\}$, we can see why the intersection of these sets of infinite numbers are actually empty. Consider a number $m$ that actually satisfies $m \in A_n$ for every $A_n$ in our collection of sets. Because $m$ is not an element of $A_{m + 1}$, no such $m$ exists and the intersection is empty.
        \item This is false. Consider the following as a counter example: Since $A_n$ is finite, consider the following sets: $A_1 = \{1\}$, $A_2 = \{2\}$, and so on. The intersection of these sets is empty because none of the sets contain numbers that are present in the other sets.
        \item This is false. Let the following sets be defined as $A = \{1\}$, $B = \{2\}$, $C = \{3\}$. If we start on the left side of the equation: $A \cap (B \cup C)$ implies $\{1\} \cap (\{2\} \cup \{3\})$ implies $\{1\} \cap \{2, 3\}$ implies $\emptyset$. From the right: $(A \cap B) \cup C$ implies $ (\{1\} \cap \{2\}) \cup \{3\}$ implies $ \emptyset \cup \{3\}$ implies $ \{\emptyset, 3\}$.
        \item  This is true. Let $x \in A$, $y \in B$, and $z \in C$. Consider the following:
        \begin{enumerate}[label=(\arabic*)]
            \item \textbf{From the left of the equation:} Inside the parenthesis, either the sets share an element and combine, or they do not share an element and are empty. This means we have two possibilities:
            \begin{enumerate}[label=\roman*.]
                \item \textbf{If the set inside the parenthesis is not empty}, then we have $x$ and $y \in A \cap B$. Then it could be the case that $x$ and $y \in C$. From this, we know 3 things can occur. Either every set contains only $y$, or $x$ and $y$.
                \item \textbf{If the set inside the parenthesis is empty}, then the resulting intersection with $C$ would be the empty set.
            \end{enumerate}
            \item \textbf{From the right of the equation:} Similar to (1), either the sets share an element and combine, and they do not share an element and are empty. This means we have two possibilities:
            \begin{enumerate}[label=\roman*.]
                \item \textbf{If the set inside the parenthesis is not empty}, then we have $y$ and $z \in B \cap C$. Then it could be the case that $y$ and $z \in A$. From this, we know 2 things can occur. Either every set contains only $y$, or $y$ and $z$.
                \item \textbf{If the set inside the parenthesis is empty}, then the resulting intersection with $A$ would be the empty set.
            \end{enumerate}
            \item In every case, either we get the empty set, or end up in a situation where every set contains just $y$ or all elements. Hence, the two sides of the equation will always equal each other. 
        \end{enumerate}
        \item This is true. Because $B \cap C$ contains all the elements in $B$ or $C$, so when the intersection is taken on either side of the equation, you will have the same elements regardless. (This statement is much shorter than the explanation given in the previous problem because I am unsure if I even a proof at all. I'm testing both waters here.) 
    \end{enumerate}
}

\newpage
\begin{exercise}
    {1.2.5}\Gls{De Morgan's Laws}. Let $A$ and $B$ be subsets of $\R$.
    \begin{enumerate}
        \item  If $x \in (A \cap B)^c$, explain why $x \in A^c \cup B^c$. This shows that $(A \cap B)^c \subseteq A^c \cup B^c$.
        \item Prove the reverse inclusion $(A \cap B)^c \supseteq A^c \cup B^c$, and conclude that $(A \cap B)^c = A^c \cup B^c$.
        \item Show $(A \cup B)^c = A^c \cap B^c$ by demonstrating inclusion both ways.
    \end{enumerate}
\end{exercise}

\sol{\hfill
    \begin{enumerate}
        \item If $x \in (A \cap B)^c$, and we know that $A^c = \{x \in \R \ \colon x \notin A\}$, then we know $x$ must cannot exist in $A^c$ and $B^c$ because $(A \cap B)^c = \{x \in \R \ \colon x \notin (A \cap B)\}$. Thus, $x$ is in either $A^c$ or $B^c$. Put another way $x \in A^c \cup B^c$. Since we have shown that an element that started in $(A \cap B)^c$ ended up in $A^c \cup B^c$, then we know $(A \cap B)^c \subseteq A^c \cup B^c$.
        \item Assume there exists a $y \in A^c \cup B^c$. Thus, it must be the case that $y \notin A$ or $y \notin B$. Hence, $y$ cannot be exist in both sets at the same time, so $y \in (A \cap B)^c$. Because we have taken an element that started in $A^c \cup B^c$ and have shown that it exists in $(A \cap B)^c$, we have proven $A^c \cup B^c \subseteq (A \cap B)^c$. 
        \item\hfill \sbseteqpf
    {Assume there exists $x \in (A \cap B)^c$, and we know that $A^c = \{x \in \R \ \colon x \notin A\}$, then we know $x$ must cannot exist in $A^c$ and $B^c$ because $(A \cap B)^c = \{x \in \R \ \colon x \notin (A \cap B)\}$. Thus, $x$ is in either $A^c$ or $B^c$. Put another way $x \in A^c \cup B^c$. Since we have shown that an element that started in $(A \cap B)^c$ ended up in $A^c \cup B^c$, then we know $(A \cap B)^c \subseteq A^c \cup B^c$.}
    {Now assume there exists a $y \in A^c \cup B^c$. Thus, it must be the case that $y \notin A$ or $y \notin B$. Hence, $y$ cannot be exist in both sets at the same time, so $y \in (A \cap B)^c$. Because we have taken an element that started in $A^c \cup B^c$ and have shown that it exists in $(A \cap B)^c$, we have proven $A^c \cup B^c \subseteq (A \cap B)^c$.}
    {Therefore, we have shown through proving both sides of the implication, that these two statements are logically equivalent. In that, all elements of $A^c \cup B^c$ are the same elements that are in $(A \cap B)^c$}
    {xgray}
    \end{enumerate}
}



\begin{exercise}
    {1.2.7}Given a function $f$ and a subset $A$ of its domain, let $f(A)$
represent the range of $f$ over the set $A$; that is, $f(a) = \{f(x) \colon x \in A\}$.
\begin{enumerate}
    \item Let $f(x) = x^2$. if $A = [0,2]$ (the closed interval $\{x \in \R \colon 0 \leq x \leq 2\}$) and $B = [1,4]$, find $f(A)$ and $f(B)$. Does $f(A\cap B) = f(A) \cap f(B)$ in this case? Does $f(A \cup B) = f(A) \cup f(B)$?
    \item Find two sets $A$ and $B$ for which $f(A\cap B) \ne f(A) \cap f(B)$.
    \item Show that, for an arbitrary function $g \ \colon \ \R \rightarrow \R$, it is always true that $g(A \cap B) \subseteq g(A) \cap g(B)$ for all sets $A,B\subseteq \R$.
    \item Form and prove a conjecture about the relationship between $g(A \cup B)$ and $g(A) \cup g(B)$ for an arbitrary function $g$.
\end{enumerate}
\end{exercise}

\sol{\hfill
    \begin{enumerate}
        \item Since $f(x) = x^2$, the intervals of $f(A)$ would be $[0,4]$ and $f(B)$ would be $[1,16]$. The interval of the intersection of $A \cap B$ is $[1,2]$. Take this through our function, we get $f(A \cap B) = [1,4]$. On the other side of the equation, we already know the intervals of $f(A)$ and $f(B)$, and the intersection of theirs would be $[1,4]$. So they do equal each other. We know $f(A \cup B)$ and $f(A) \cup f(B)$ will be equivalent because $f(A \cup B)$ has an interval of $[0,16]$, and $f(A) \cup f(B)$ also has an interval of $[0,16]$ because taking the union of $[0,4] \cup [1,16]$ is $[0,16]$.
        \item Two sets could be $A = [5,6]$ and $B = [0,0]$. Because the sets have nothing in common even after taking their function, they do not equal each other. 
        \item \hfill \expf{Let $x \in g(A \cap B)$. Using the definition of function, we know there exists a $y \in A \cap B$ to which that $y$ is mapped to as $g(y) = x$. From the definition of intersection, we know $y \in A$ and $y \in B$ such that $x = g(y) \in g(A)$ and $x = g(y) \in g(B)$ because $y \in A \cap B$. Putting it together, we have $x \in g(A) \cap g(B)$ thus proving $g(A \cap B) \subseteq g(A) \cap g(B)$}
        \item Conjecture: For any function $g$ defined as $g \ \colon \R \rightarrow \R$ and for any subsets $A,B \subseteq \R$, the following holds: $$g(A \cup B) = g(A) \cup g(B)$$ 
        \sbseteqpf
        {Take any element $x \in g(A \cap B)$. By definition of function, we know there exists some $y \in A \cup B$ such that $g(x) = y$. From the definition of union, we know $y \in A$ or $y \in B$ such that $x = g(y) \in g(A)$ or $x = g(y) \in g(B)$ or both. Putting it together, we have $x \in g(A) \cup g(B)$ thus proving $g(A \cup B) \subseteq g(A) \cup g(B)$.}
        {Take any element $p \in g(A) \cap g(B)$. By definition of union, we know $p$ is either in $g(A)$ or $g(B)$ or both. From the definition of function, we know that if $p \in g(A)$ or $p \in g(B)$ then there exists some $q \in A$ or $q \in B$ such that $g(q) = p$. Putting it together, we have $q \in A \cup B$. Moreover, this means $p = g(x) \in g(A \cup B)$. And since $p \in g(A) \cup g(B)$ implies $p \in g(A \cup B)$, we know $g(A) \cup g(B) \subseteq g(A \cup B)$.}
        {Therefore, since we have proven that both expressions are functions of each other, we have proved that they are equal.}
        {xgray}
    \end{enumerate}
}

\begin{exercise}
    {1.2.8}Given a function $f \ \colon A \rightarrow B$ can be defined as either \gls{one-to-one} or \gls{onto}, give an example of each or state that the request is impossible:
    \begin{enumerate}
        \item $f \ \colon \N \rightarrow \N$ that is 1-1 but not onto.
        \item $f \ \colon \N \rightarrow \N$ that is onto but not 1-1.
        \item $f \ \colon \N \rightarrow \Z$ that is 1-1 and onto.
    \end{enumerate}
\end{exercise}

\sol{\hfill
    \begin{enumerate}
        \item The function $f(a) + 1$ is 1-1 because when \begin{align*}
            f(a_1) &= f(a_2)\\
            a_1 + 1 &= a_2 + 1 \\
            a_1 &= a_2
        \end{align*}
        However, the function is not onto because the entire co-domain is not covered. That being 1.
        \item \textit{We need to find a function that will cover every entry in the co-domain, while also avoiding a scenario where $a_1 = a_2$...} Consider the function, 
        $$f(a) =
        \begin{cases}
            a &\text{ if \textit{a} is odd,}\\
            a - 1 &\text{ if \textit{a} is even}
        \end{cases}$$
        This function is onto because every natural number is covered, but it is not 1-1 because $a_1 \ne a_2 - 1$.
        \item This request is not possible. There is no way to map every natural number to every integer because we are simply missing 0! (Not 0 factorial, we do have the number 1, I just mean the number 0 in a exclamatory sense.)
    \end{enumerate}
}
\setcounter{section}{2}
\section{Axiom of Completeness}

Think about $\Q$ and $\R$.
    \begin{itemize}
        \item Both are fields.
        \begin{itemize}
            \item Both have $+,-,\times,\div$ operations.
        \end{itemize}
        \item Both are totally ordered
        \begin{itemize}
            \item $a < b$,
            \item $a > b$, 
            \item or $a = b$
        \end{itemize}
        \item $\R$ is complete. $\Q$ is not.
    \end{itemize}
    
\begin{axiom}
    {Completeness}Every nonempty set of real numbers that is \textbf{\gls{bounded}} has a \textbf{\gls{least upper bound}}.
\end{axiom}

\textbf{Note that \glspl{upper bound} are not unique!} For example, consider the line, $A$, from $0$ to $1$. There are infinitely many upper bounds past $1$ because $A$ is bounded. \\

We often call the least upper bound the \textit{supremum} of a set. Example: \\

Imagine a number line from $(1,8)$. Note that parenthesis mean $<$ and not $\leq$. Hence, the supremum is $8$. Wrote simply as $\sup A$. \\

\begin{example}
    {Supremum}Consider a set, $B = [-5,-2] \cup (3,6) \cup \{13\}$. What is the supremum?
\end{example}

\sol{
    $\sup B = 13$
}

At the other end of the set, we have the following:
\begin{itemize}
    \item lower bounds,
    \item greatest lower bound
    \item often called infimum.
\end{itemize}

The infimum of the previous example would be $\inf B = -5$.

\begin{example}
    {}Consider the set, $\C = \{\frac{1}{n} \ \colon \ n \in \N\}$. What is the supremum and the infimum?
\end{example}

\lesol{
    $\sup \C = 1$, $\inf \C = 0$.
}

\begin{example}
    Let $A \subseteq \R$ be nonempty and bounded above, and let $c \in \R$. Define the set $c + A$ by $$c + A = \{c+a \ \colon \ a \in A\}$$ Then $\sup(c + A) = c + \sup A$.
\end{example}

\lesol{
To properly verify this we focus separately on each part of Definition 1.3.2. Setting $s = \sup A$, we see that $a \leq s$ for all $a \in A$, which implies $c + a \leq c + s$ for all $a \in A$. Thus, $c + s$ is an upper bound for $c + A$ and condition (i) is verified. For (ii), let $b$ be an arbitrary upper bound for $c + A$; i.e., $c + a \leq b$ for all $a \in A$. This is equivalent to $a \leq b - c$ for all $a \in A$, from which we conclude that $b - c$ is an upper bound for $A$. Because $s$ is the least upper bound of $A$, $s \leq b - c$, which can be rewritten as $c + s \leq b$. This verifies part (ii) of Definition 1.3.2, and we conclude $\sup (c + A) = c + \sup A$.
}

Why do we need to include infimum and supremum? Don't we have the max and min of a set already? Well, what exactly do we mean by the \gls{maximum value} of a set? \\

We say $m \in \R$ is the \textit{maximum} of $A$ if $m \in A$ and for all $x\in A$, $x \leq m$. Note that some sets have a maximum and some sets do not. You cannot refer to a maximum without first knowing it exists. This is the same with minimums. Suprema do not exist if the set is bounded.

\begin{lemma}
    {1.3.8}Assume $s$ is an \gls{upper bound} for a set $A \subseteq \R$. Then, $s$ is the supremum of $A$ if and only if for every $\epsilon > 0$, there exists $x\in A$ such that $s - \epsilon < x$.
\end{lemma}

This lemma allows us to take any positive number and take a ``step back.'' In essence, you can verify something as an upper bound if you continuously back up over and over until you cannot back up any longer.

\iffpf
{
Assume $s = \sup A$. Let $\epsilon > 0$. Suppose there are no elements $x$ of $A$ such that $s - \epsilon < x$. Then $s - \epsilon$ would be an upper bound. This contradicts that $s$ is the least upper bound. Therefore, there must exist an element $x \in A$ such that $s - \epsilon < x$. 
}{
Assume for every $\epsilon > 0$, there exists $x \in A$ such that $s - \epsilon < x$. Let $t$ be an upper bound of $A$. Suppose $t < s$. Consider $\epsilon_0 = s - t > 0$. By our assumption, there exists $x \in A$ such that $s - \epsilon_0 < x$. So, $t < x$. This contradicts that $t$ is an upper bound of $A$. So, $t \geq s$. Thus, $s$ is the least upper bound
}{
Therefore, by proving both the right and left implication, we have shown the statement to be true.
}{xpurple}
Analogous statement about infimums: Assume $z$ is a lower bound of a set $A \subseteq \R$. Then $z = \inf A \iff $ for all $ \epsilon >0$, there exists $y \in A$ such that $y < z + \epsilon$.

\subsection*{Exercises}

\begin{exercise}
    {1.3.4}Let $A_1,A_2,A_3\dots$ be a collection of nonempty sets each of which is bounded above. 
    \begin{enumerate}
        \item Find a formula for $\sup(A_1 \cup A_2)$. Extend this to $\sup (\bigcup^n_{k=1}A_k)$.
        \item Consider $\sup(\bigcup^\infty_{k=1} A_k)$. Does the formula in (a) extend to the infinite case?
    \end{enumerate}
\end{exercise}

\sol{\hfill
\begin{enumerate}
    \item Let $A_1$ and $A_2$ be nonempty sets, each bounded above. To find the largest of the two suprema, we can use the following: $\sup(A_1 \cap A_2) = \max\{\sup A_1, \sup A_2\}$. If we extend this notion to $\sup (\bigcup^n_{k=1}A_k)$, we can use the same idea from before and write it as $\sup (\bigcup^n_{k=1}A_k) = \max\{\sup A_1, \sup A_2, \dots, \sup A_n\}$.
    \item The formula does not extend to the infinite case. Consider the counterexample $\bigcup^\infty_{k=1} A_k$ where $A_k := [k, k+1]$. Even though these sets are bounded above, when we take the union of them, we approach infinity, which is not bounded: $\bigcup^\infty_{k=1} A_k = [1,2] \cup [2,3] \cup \dots = [1, \infty)$.
\end{enumerate}
}

\begin{exercise}
    {1.3.5}As in Example 1.3.7, let $A \subseteq \R$ be nonempty and bounded above, and let $c \in \R$. This time define the set $cA = \{ca \ \colon a \in A\}$.
    \begin{enumerate}
        \item If $c \geq 0$, show that $\sup(cA) = c\sup A$.
        \item Postulate a similar type of statement for $\sup(cA)$ for the case $c < 0$.
    \end{enumerate}
\end{exercise}

\sol{
\begin{enumerate}
    \item Let $A \subseteq \R$ be nonempty and bounded above. Define the set $cA := \{ca \ \colon a \in A\}$. From the axiom of completeness, because $A$ is bounded above, we know there is a least upper bound, $s = \sup A$. Following from Example 1.3.7, we see that $a \leq s$ for all $a \in A$ which implies $ca \leq cs$ for all $a \in A$. Thus, $cs$ is an upper bound for $cA$, and the first condition of Definition 1.3.2 is satisfied. For the second condition, we need to look at both $c = 0$ and $c > 0$ to avoid dividing by zero. So, we have two cases:
    \begin{itemize}
        \item \textbf{$c = 0$:} If $c = 0$, then $cA = \{0 \colon \ a \in A\} = \{0\}$. Since the only element in $cA$ is $0$, $\sup(cA) = 0$. Similarly, because $c = 0$, $c \sup A = 0 \cdot \sup A = 0$. Therefore, $\sup(cA) = c \sup(A)$.
        \item \textbf{$c > 0$:} Let $b$ be an arbitrary upper bound for $cA$ and $c > 0$. In other words, $ca \leq b$ for all $a \in A$. This is equivalent to $a \leq b / c$ where $c \ne 0$, from which we can see that $b / c$ is an upper bound for $A$. Because $s$ is the least upper bound of $A$, $s \leq b / c$, which can be rewritten as $cs \leq b$. This verifies the second part of Definition 1.3.2, and we conclude $\sup(cA) = c \sup A$.
    \end{itemize}
    \item Postulate: If $c < 0$, then $\sup(cA) = c \inf (A)$. \\
    % \expf{Assume $c < 0$, Multiplying $a \leq s$ by $c$ reverses the inequality, so $ca \geq cs$ for all $a \in A$. Therefore, $cs$ is a lower bound of for $cA$. To verify this as the lowest upper bound, let $m  = \inf A$. Using the same logic as before, for all $a \in A$, $m \leq a$, so $cm \geq ca$. This shows that $cm$ is the least upper bound of $cA$, and thus $\sup(cA) = cm = c \inf A$.}
\end{enumerate}
}

\begin{exercise}
    {1.3.8}Compute, without proofs, the suprema and infima (if they exist) of the following sets:
\begin{enumerate}
    \item $\left\{\frac{m}{n} : m, n \in \mathbb{N} \text{ with } m < n\right\}$.
    \item $\left\{\frac{(-1)^m}{n} : m, n \in \mathbb{N}\right\}$.
    \item $\left\{\frac{n}{3n+ 1} : n \in \mathbb{N}\right\}$.
    \item $\left\{\frac{m}{m+ n} : m, n \in \mathbb{N}\right\}$.
\end{enumerate}
\end{exercise}

\sol{To avoid writing out every set definition, I am going to denote each set as $A_n$ where $n$ corresponds to the numerical value of the list from (a) - (d).
\begin{enumerate}
    \item $\sup A_1 = 1$, $\inf A_1 = 0$
    \item $\sup A_2 = 1$, $\inf A_2 = -1$
    \item $\sup A_3 = \frac{1}{3}$, $\inf A_3 = \frac{1}{4}$
    \item $\sup A_4 = 1$, $\inf A_3 = 0$
\end{enumerate}
}

\section{Consequences of Completeness}

\begin{ntheorem}
    {Nested Interval Property}For each $n \in \N$, assume we are given a closed interval $I_n = (a_n, b_n)$. Assume $I_n$ contains $I_{n+1}$. This results in a nested sequence of intervals. $$I_1 \supseteq I_2 \supseteq I_3 \supseteq I_4\dots$$ Then, $\bigcap^\infty_{n=1} \ne \emptyset$.
\end{ntheorem}

\textbf{tl;dr} there has to be something that is common to all of the sets.

\tpf{
    Notice that the sequence, $a_1,a_2,a_3,\dots$ is increasing. In other words, for each $n \in \N$, since $I_n \supset I_{n + 1}$ we have $a_n \leq a_{n+1}$. If we consider the set $A = \{a_n \ \colon \ n\in \N\}$. The element $b_1$ is an upper bound of $A$. (Note that $b_1$ and $a_1$ corresponds to the endpoints of the first set, $I_1$. Think of this as a tornado looking structure where the larger the $I_n$, the smaller the number line.) For each $n \in \N$, $a_n \leq b_n \leq b_1$. \\

    Since $A$ has an upper bound, it must have a least upper bound. Hence, let $\alpha = \sup A$. We claim that $\alpha \in \bigcap^\infty_{n=1} I_n$. We said $b_1$ was an upper bound. In fact, every $b_n$ is an upper bound of $A$. Choose any $n,m \in \N$. We want to show that $a_n \leq b_m$. Consider the following cases: \\

    \textbf{Case 1:} If $n < m$, then $a_n \leq a_m \leq b_m$. (Think: two number lines stacked on top of each other. The top number line is larger, call it $I_n$ and it has $a_n$ and $b_n$ as endpoints. Consider a contained line ($I_n \supseteq I_m$) that is smaller, and has endpoints $a_m$ and $b_m$.) \\

    \textbf{Case 2:} If $n > m$, then $a_n \leq b_n \leq b_m$. So every $b_n$ is an upper bound of $A$.\\

    Hence, 
    \begin{itemize}
        \item Because $\alpha = \sup A$, we have $\alpha \geq a_n$. 
        \item Since $b_n$ is an upper bound of $A$, we have $\alpha \leq b_n$.
    \end{itemize}
    so, $\alpha \in [a_n, b_n] = I_n$. Thus, $\alpha \in \bigcap^\infty_{n=1} I_n.$
}

Nested, closed, Bounded Intervals $\Rightarrow$ non-empty intersection.

\begin{ntheorem}
    {Archimedean Principle} 
    \begin{enumerate}
        \item Given any number $x \in R$, there exists an $n \in N$ satisfying $n > x$.
        \item Given any real number $y > 0$, there exists an $n \in N$ satisfying $1/n < y$.
    \end{enumerate}
\end{ntheorem}

\tpf{
    \begin{enumerate}
        \item If $\N$ was bounded, then we can let $s \in \N = \sup \N$. However, we know that there is always a higher number (e.g., $n+1$) for any $n \in \N$ that is given. Thus, by contradiction, there must exist $n \geq x$.
        \item For any $x > 0$, there exists $n \in \N$ such that $\frac{1}{n} < x$. 
    \end{enumerate}
}

\begin{theorem}
    Density of $\Q$ in $\R$. For any $a,b \in \R$ with $a < b$, there exists $q \in \Q$ such that $a < q < b$.
\end{theorem}

\tpf{
Since $ b - a > 0$, there exists $n \in \N$ such that $\frac{1}{n} < b - a$. From the \namrefthm{Archimedean Principle}, since $a\times n \in \R$, there exists $m \in \N$ such that $a\times n < m$. Let $m$ be there smallest such natural numbers (by the well ordered principle). Since $m$ is the smallest such natural number, it follows that $m - 1 \leq a\times n < m$. We then see that $a < \frac{m}{n}$. Now, we need to find some $\frac{m}{n} < b$. \begin{align*}
    m - 1 &\leq a\times n \\
    m &\leq a \times n + 1 \\
    \frac{m}{n} &\leq a + \frac{1}{n} \\
    \frac{m}{n} &< a + (b - a) \\
    \frac{m}{n} &< b
\end{align*} We now have that $a < \frac{m}{n} < b$ so $ \frac{m}{n}$ is a rational number in $(a,b)$
}

\begin{exercise}
    {1.4.1}Recall that $\mathbb{I}$ stands for the set of irrational numbers. 
    \begin{enumerate}
        \item Show that if $a,b \in \Q$, then $ab$ and $a + b$ are elements of $\Q$ as well.
        \item Show that if $a \in \Q$ and $t \in \mathbb{I}$, then $a + t \in \mathbb{I}$ and $at \in \mathbb{I}$ as long as $a \ne 0$.
        \item Part (a) can be summarized by saying that $\Q$ is closed under addition and multiplication. Is $\mathbb{I}$ closed under addition and multiplication? Given two irrational numbers $s$ and $t$, what can we say about $s + t$ and $st$? In other words, are there two irrational numbers that can be added and multiplied such that you get a number $x$ such that $x \notin \mathbb{I}$.
    \end{enumerate}
\end{exercise}

\sol{
\begin{enumerate}
    \item Let $a,b \in \Q$. This means there exists some $p,q,a,b \in \Z$ such that $$a = \frac{p}{q}$$ and $$b = \frac{a}{b}$$ where $q,b \ne 0$. The product of these numbers is $$ab = \frac{p}{q} \cdot \frac{a}{b} = \frac{pa}{qb}.$$ Since $pa,qb \in \Z$, $ab \in \Q$. The sum of these numbers is $$a + b = \frac{p}{q} + \frac{a}{b} = \frac{pb + aq}{qb}.$$ Since $pb + aq, qb \in \Z$, $a + b \in \Q$. 
    \item Let $a \in \Q$ and $t \in \I$. Assume, for contradiction, that $a + t \in \Q$. This would imply $t = (a + t) - a$ (because we can subtract $t + a$ from the original equation and rearrange terms). Since $a + t, a \in \Q$ their sum would be rational because the rational numbers are closed under addition. However, that would contradict the assumption that $t \in \I$. Hence, $a + t \in \I$.   
    \item For $\I$, it is not closed under addition and multiplication. Consider the following counterexample: $\sqrt{2} + (-\sqrt{2}) = 0$ which is not in the irrationals. For multiplication, consider $\sqrt{2} \cdot \sqrt{2} = 2$, which is also not in the irrationals. 
\end{enumerate}
}

\section{Cardinality}

Two sets have the same \textit{cardinality} if there exists a bijection between them. Thus, the natural numbers, the integers, and the rational numbers have the same cardinality. A set is \textit{countably infinite} if it has the same cardinality as $\N$. (If it can be put into one-to-one correspondence with $\N$.) A set is \textit{countable} if it is countably infinite or finite. \\

\setcounter{BoxCounter}{5}
\begin{theorem}
    $\R$ is not countable.
\end{theorem}

\tpf{%
    \textit{1} (most common) \\
    Suppose $\R$ is countable. Then we can list them all, or we can enumerate them. $\R = \{x_1,x_2,x_3,x_4,\dots\}$. We can write the decimal expansion of each of these. Consider the following table:
\begin{center}
\begin{tabular}{lllllll}
$x_1 =$ & $\boxed{a_{10}}$ & $a_{11}$ & $a_{12}$ & $a_{13}$ & $a_{14} $ & \dots \\ \hline
$x_2 =$ & $a_{20} $ & $\boxed{a_{21}}$ & $a_{22} $ & $a_{23} $ & $a_{24} $ & \dots \\ \hline
$x_3 =$ & $a_{30} $ & $a_{31} $ & $\boxed{a_{32}}$ & $a_{33} $ & $a_{34} $ & \dots \\ \hline
$x_4 =$ & $a_{40} $ & $a_{41} $ & $a_{42} $ & $\boxed{a_{43}}$ & $a_{44} $ & \dots \\ \hline
$x_5 =$ & $a_{50} $ & $a_{51} $ & $a_{52} $ & $a_{53} $ & $\boxed{a_{54}}$ & \dots \\ \hline
$x_6 =$ & $a_{60} $ & $a_{61} $ & $a_{62} $ & $a_{63} $ & $a_{64} $ & \dots
\end{tabular}
\end{center}

We will now construct a number that is not in this list. Focus on diagonal entries. For each $n \in \N$, let $b_n$ be a digit that is different fron $a_{nn}$. Now consider the number $y = 0.b_1b_2b_3b_4b_5\dots$. This number $y$ is not in our list. So our list did not include all of $\R$. Avoid repeating $9$s.
}

\tpf{
\textit{2} (uses nested interval theorem) \\
Suppose $\R$ is countable. Then we can enumerate $\R$ $\R = \{x_1,x_2,x_3,\dots\}$. Let $I_1$ be any closed interval that does not contain $x_1$. Next, we will find another closed interval $I_2$ that:\begin{itemize}
    \item $I_2 \subseteq I_1$ 
    \item $x_2 \notin I_2$
\end{itemize} Continue in this fashion creating a sequence of nested closed intervals: $I_1\supseteq I_2 \supseteq I_3 \supseteq \dots$ such that for all $k \in \N$, $x_k \notin I_k$. Now consider: $$\bigcap^\infty_{n=1} I_n$$
\begin{itemize}
    \item For each $k \in \N$, since $x_k \notin I_k$, we see $x_k \notin \bigcap^\infty_{n=1} I_n$.
    \item By the nested interval theorem, there exists $x \in \R$ such that $x \in \bigcap^\infty_{n=1} I_n$. So $x$ is a real number that is not included in our list.
\end{itemize}
}

\begin{theorem}
    A countable collection of finite sets is \textit{countable}.
\end{theorem}

\begin{theorem}
    \begin{enumerate}[label=(\roman*)]
        \item The union of two countable sets is \textit{countable}.
        \item A countable union of countable sets is \textit{countable}.
    \end{enumerate}
\end{theorem}

From Theorem \ref{thm:1.5.6}, we know that $\R$ is uncountable, but what about $(0,1)$? It does have the same cardinality of $\R$ because we can make a one-to-one and onto function between both the sets. Similarly, $(a,b)$ also has the same cardinality. What about $[a,b]$? 

\textbf{Recap:} $\N$ is countable, and $\R$ is uncountable and has a different cardinality than $\N$. Thus, the question is, do all uncountable sets have the same cardinality as $\R$? The answer is \textbf{no}.
    
\begin{ntheorem}
    {Canter's Theorem} For any set $A$, there does not exist an onto map from $A$ into $\mathcal{P}$.
\end{ntheorem}

\tpf{
    Suppose there exists an onto function, $f \ \colon \ A \rightarrow \mathcal{P}(A)$. So each $a \in A$ is mapped to an element $f(a) \in \mathcal{P}(A)$. Then, $f(a) \subseteq A$. We are going to construct an element of $\mathcal{P}(A)$ which is not mapped to by $f$. 

    Consider $B = \{a \in A \ \colon \ a \notin f(a)\}$. Since $f$ is onto there exists $a' \in A$ such that $B = f(a')$. Thus, there are two cases to consider:

    \begin{itemize}
        \item \textbf{Case 1:} If $a' \in B = f(a')$, then $a' \notin B$.
        \item \textbf{Case 2:} If $a' \notin B = f(a')$, then $a' \in B$.
    \end{itemize}

    As evidenced, both cases lead to contradictions, so $B$ is not the image of any $a \in A$. Therefore $f$ is not onto.
}

\begin{example}
    {Set and Power Set Matching}$A = \{a,b,c\}$. 
\end{example}

\sol{
$\mathcal{P}(A) = \emptyset, \{a\}, \{b\}, \{c\}, \{a,b\}, \{a,c\}, \{b,c\}, \{a,b,c\}$. Note that you can map $\{a\}$, $\{b\}$, $\{c\}$, to elements such as $\emptyset$, $\{a,b\}$, $\{a,b,c\}$, but there are still more elements that are left unmapped. We can extrapolate from our proof a set $B$ such that $B = \{a,c\}$ because those elements are not mapped to.
}

All of this is to show $\mathcal{P}(\R)$ has a larger cardinality than $\R$. Then $\mathcal{P}(\mathcal{P}(\R))$ has a larger cardinality than $\mathcal{P}(\R)$.

% \chapter{Power Set} \label{Power Set}
% \vspace*{-0.25in}
% \input{Chapters/ch_6}

% \chapter{Equivalency} \label{Equivalency}
% \vspace*{-0.25in}
% \input{Chapters/ch_7}

% \chapter{Functions \& Bijections}\label{Functions}
% \vspace*{-0.25in}
% \input{Chapters/ch_8}

% \chapter{Recursion \& The Whole Numbers}\label{Recursion}
% \vspace*{-0.25in}
% \input{Chapters/ch_9}

% \chapter{The Integers} \label{Integers}
% \vspace*{-0.25in}
% \input{Chapters/ch_10}

% \chapter{Cardinality} \label{Cardinality}
% \vspace*{-0.25in}
% \input{Chapters/ch_11}

\printglossaries
%-%-%-%-%-%-%-%-%-%-%-%-%-%-%-%-%-%-%-%-%-%-%-%-%-%-%-%-%-%-%-%-%-%-%-%-%-%-%

\end{document}
%-%-%-%-%-%-%-%-%-%-%-%-%-%-%-%-%-%-%-%-%-%-%-%-%-%-%-%-%-%-%-%-%-%-%-%-%-%-%

%-%-%-%-%-%-%-%-%-%-%-%-%-%-%-%-%-%-%-%-%-%-%-%-%-%-%-%-%-%-%-%-%-%-%-%-%-%-%